% 太阳系外行星（综述）
% license CCBYSA3
% type Wiki

本文根据 CC-BY-SA 协议转载翻译自维基百科\href{https://en.wikipedia.org/wiki/Exoplanet}{相关文章}。

\begin{figure}[ht]
\centering
\includegraphics[width=8cm]{./figures/c4a53fef76fc3e3b.png}
\caption{由 W. M. 凯克天文台在七年间拍摄的 HR 8799 系统中的四颗系外行星影像，行星的运动由年度观测结果插值显示} \label{fig_Exopla_2}
\end{figure}
\begin{figure}[ht]
\centering
\includegraphics[width=8cm]{./figures/fa92c1fc35546edf.png}
\caption{开普勒-37 系统中系外行星的尺寸与水星、火星和地球的大小对比} \label{fig_Exopla_1}
\end{figure}
系外行星（exoplanet，或称太阳系外行星）是指位于太阳系之外的行星。首个得到确认的系外行星发现于 1992 年，环绕一颗脉冲星运行；而围绕主序星的首次确认发现发生在 1995 年。另有一颗行星最早于 1988 年被探测到，并在 2003 年得到确认。2016 年，人们认识到关于系外行星的最早可能证据甚至可以追溯到 1917 年。截止至 2025 年 12 月 4 日，已确认的系外行星共有 6,053 颗，分布在 4,510 个行星系统中，其中 1,022 个系统包含多于一颗行星$^{[1][2]}$。

探测系外行星的方法有多种，其中凌日测光法和多普勒光谱法发现的数量最多。但这些方法存在明显的观测偏差，更容易发现靠近母恒星的行星；因此，已探测到的系外行星中约有 85\% 位于潮汐锁定区之内$^{[3]}$。据估计，大约每 5 颗类太阳恒星中就有 1 颗拥有一颗位于宜居带内的“类地尺寸”行星$^{[a][b][c][4][5]}$。若假设银河系中约有 2,000 亿颗恒星$^{[d]}$，则可以推测银河系内可能存在约 110 亿颗潜在宜居的类地尺寸行星；若将数量众多的红矮星系统也纳入考虑，这一数字可上升至约 400 亿颗$^{[6]}$。

目前已知质量最小的系外行星是 Draugr，其质量约为月球的两倍。NASA 系外行星档案中列出的质量最大的系外行星是 HR 2562 b$^{[7][8][9]}$，其质量约为木星的 30 倍。然而，根据某些以氘核聚变为界限的行星定义$^{[10]}$，它的质量可能过大，更应被归类为褐矮星。已知系外行星的轨道周期跨度极大，从不足一小时（最靠近母恒星者）到数千年不等。有些系外行星距离母恒星极远，以至于难以判断它们是否仍与恒星保持引力束缚。

距离地球最近的系外行星位于约 4.2 光年（1.3 秒差距）之外，环绕着比邻星运行，比邻星是距离太阳最近的恒星$^{[11]}$。在另一极端，人们也发现了关于河外行星的证据，即位于其他星系中的系外行星$^{[12][13]}$。

系外行星的发现极大地激发了人们对地外生命搜索的兴趣，尤其关注那些位于恒星宜居带（有时称为“金发姑娘区”）内的行星，因为在那里行星表面可能存在液态水——这是我们所知生命的必要条件。然而，对行星宜居性的研究还需要综合考虑诸多其他因素，以评估一颗行星是否适合孕育生命$^{[14]}$。詹姆斯·韦布空间望远镜（JWST）与地基及其他空间望远镜协同工作，有望在系外行星的成分、环境条件及宜居性等方面提供更深入的认识$^{[15]}$。

流浪行星是指不隶属于任何行星系统的行星。这类天体通常被视为独立于行星之外的一个类别，尤其当它们是气体巨行星时，往往被归类为次级褐矮星$^{[16]}$。银河系中流浪行星的数量可能高达数十亿颗$^{[17][18]}$。
\subsection{定义}
\subsubsection{IAU（国际天文学联合会）}
国际天文学联合会（IAU）所使用的“行星”这一术语的官方定义仅适用于太阳系，因此并不直接适用于系外行星$^{[19][20]}$。IAU 系外行星工作组（Working Group on Extrasolar Planets）于 2001 年发布了一份立场声明，其中给出了“行星”的工作性定义，并于 2003 年进行了修订$^{[21]}$。当时，系外行星被定义为满足以下条件的天体：
\begin{itemize}
\item 真质量低于氘发生热核聚变的临界质量（对于太阳金属丰度的天体，目前计算值约为 13 个木星质量），且环绕恒星或恒星遗骸运行的天体，被定义为“行星”（不论其形成机制如何）。判定一个系外天体是否为行星所需的最小质量或尺寸，应与太阳系中采用的标准相同。
\item 真质量高于氘发生热核聚变临界质量的亚恒星天体，被定义为“褐矮星”，不论其形成方式或所在位置如何。
\item 在年轻疏散星团中发现的、质量低于氘热核聚变临界质量且不受恒星束缚而自由漂浮的天体，不被视为“行星”，而应归类为“次级褐矮星”（或其他更合适的名称）。
\end{itemize}
这一工作性定义随后在 2018 年 8 月由 IAU 的 F2 委员会（系外行星与太阳系委员会）进行了修订$^{[22][23]}$。目前，IAU 对系外行星的官方工作定义如下：
\begin{itemize}
\item 真质量低于氘发生热核聚变的临界质量（对于太阳金属丰度天体，目前计算值约为 13 个木星质量），并且环绕恒星、褐矮星或恒星遗骸运行，同时其与中心天体的质量比低于 L4/L5 拉格朗日点不稳定性条件（$M/M_{\text{central}} < \frac{2}{25+\sqrt{621}}$）的天体，被定义为“行星”（不论其形成机制如何）。
\item 判定一个系外天体是否为行星所需的最小质量或尺寸，应与太阳系中采用的标准相同。
\end{itemize}
\subsubsection{替代方案}
IAU 的工作性定义并非始终被采用。另一种替代观点认为，应当依据**形成机制**来区分行星与褐矮星。普遍认为，巨行星通过“核吸积”形成，这一过程有时可能产生质量超过氘聚变阈值的行星$^{[24][25][10]}$；此类大质量行星可能已经被观测到$^{[26]}$。褐矮星则类似恒星，是由气体云的直接引力坍缩形成的，而这种形成机制同样可以产生质量低于 13 个木星质量、甚至低至 1 个木星质量的天体$^{[27]}$。
在这一质量范围内、以数百或数千天文单位（AU）的宽分离轨道绕恒星运行、且恒星/天体质量比很大的对象，很可能是以褐矮星方式形成的；其大气成分预计更接近其宿主恒星，而非通过吸积形成的行星（后者通常富含较重元素）。截至 2014 年 4 月，大多数被直接成像的“行星”质量较大且轨道很宽，因此很可能代表了褐矮星形成机制的低质量端$^{[28]}$。有研究提出，质量高于 10 个木星质量的天体是通过引力不稳定形成的，不应被视为行星$^{[29]}$。

此外，“13 个木星质量”的分界线并不存在精确的物理意义。氘聚变在某些质量低于该阈值的天体中同样可能发生$^{[10]}$，且发生的氘聚变量在一定程度上取决于天体的化学组成$^{[30]}$。2011 年，《系外行星百科全书》收录了质量高达 25 个木星质量的天体，并指出：“在观测到的质量谱中，13 个木星质量附近并不存在特殊特征，这进一步强化了放弃该质量上限的选择”$^{[31]}$。截至 2016 年，基于对质量—密度关系的研究，该上限被提高到 60 个木星质量$^{[32][33]}$。
Exoplanet Data Explorer 收录了质量最高达 24 个木星质量的天体，并附注：“IAU 工作组所采用的 13 个木星质量区分标准在物理上缺乏动机（尤其是对具有岩质核心的行星而言），且在观测上由于 $\sin i$ 不确定性而存在问题”$^{[34]}$。NASA 系外行星档案库则收录质量（或最小质量）不超过 30 个木星质量的天体$^{[35]}$。另一个区分行星与褐矮星的判据（不同于氘聚变、形成过程或轨道位置）是：其核心压力是由库仑压力主导，还是由电子简并压主导，两者的分界大约在 5 个木星质量附近$^{[36][37]}$。
\subsubsection{确认}
在 NASA 的系外行星档案库中，一颗系外行星的确认标准是：要么“不同的观测技术揭示了只能由行星解释的特征”$^{[38]}$，要么通过分析性技术加以确认$^{[2]}$。
对于《系外行星百科全书》而言，“若某一天体在经认可的论文或专业会议上被明确无歧义地主张为行星，则被视为已确认”$^{[39]}$。
\subsection{命名法}
\begin{figure}[ht]
\centering
\includegraphics[width=6cm]{./figures/a30e59faab567d01.png}
\caption{系外行星 HIP 65426b 是围绕恒星 HIP 65426 发现的第一颗行星$^{[40]}$。} \label{fig_Exopla_3}
\end{figure}
系外行星的命名惯例是国际天文学联合会（IAU）所采用的多恒星系统命名体系的延伸。对于绕单一恒星运行的系外行星，其 IAU 指定名称由母恒星的指定名或专名加上一个小写字母构成$^{[41]}$。字母按照围绕母恒星发现行星的先后顺序依次分配，因此在一个行星系统中，最先发现的行星被标记为“b”（母恒星本身视为“a”），随后发现的行星依次使用后续字母。如果同一系统中的多颗行星在同一时间被发现，则距离恒星最近的行星获得下一个字母，其余行星按轨道尺度顺序依次命名。IAU 还制定了一套临时认可的标准，用于处理环双星行星的命名。少数系外行星已被 IAU 正式赋予专名，此外还存在其他命名体系$^{[citation\ needed]}$。
\subsection{探测历史}
\begin{figure}[ht]
\centering
\includegraphics[width=8cm]{./figures/d13a18693bd77d1a.png}
\caption{截至 2022 年的现有及未来系外行星任务的 NASA 示意图。} \label{fig_Exopla_4}
\end{figure}
几个世纪以来，科学家、哲学家以及科幻作家一直怀疑系外行星的存在，但始终无法确定它们是否真实存在、数量是否普遍，或是否与太阳系行星相似。19 世纪曾提出过多种探测主张，但均被天文学家否定（有待引文）。

1917 年，人们记录到了围绕范·马南 2（Van Maanen 2）运行的一颗可能的系外行星的最早证据，但直到 2016 年才被认识为系外行星的迹象。天文学家沃尔特·西德尼·亚当斯（Walter Sydney Adams）使用威尔逊山天文台的 60 英寸望远镜获得了该恒星的光谱，并将其解释为 F 型主序星的光谱。后来，在对白矮星异常成分的研究中，这一光谱被重新审视。如今认为，这种光谱可能是附近一颗系外行星在恒星引力作用下被粉碎，其残余尘埃随后落到恒星表面所造成的结果。

20 世纪中期还出现过多起其他发现主张，涉及天鹅座 61、拉兰德 21185 以及巴纳德星，但这些主张直到 20 世纪 70 年代中后期才被否定（见下文“被否定的发现”）。1988 年，又有一次被认为是系外行星的疑似科学探测。随后不久，1992 年，亚历山大·沃尔什昌（Aleksander Wolszczan）和戴尔·弗雷尔（Dale Frail）宣布在毫秒脉冲星 PSR B1257+12 周围发现了两颗质量与类地行星相当的行星，这是目前被普遍接受的首次系外行星发现。1995 年，天文学家确认了第一颗围绕主序星运行的系外行星——一颗巨行星，以四天的周期绕附近恒星 51 飞马座运行。此后，一些系外行星被望远镜直接成像，但绝大多数仍是通过凌日法和径向速度法等间接方法探测到的。
\begin{figure}[ht]
\centering
\includegraphics[width=8cm]{./figures/85a4bfec6f369e50.png}
\caption{截至 2024 年 9 月，每年探测到的系外行星数量 $^{[46]}$} \label{fig_Exopla_5}
\end{figure}
2018 年 2 月，研究人员利用**钱德拉 X 射线天文台**并结合一种称为**微引力透镜**的行星探测技术，在一个遥远的星系中发现了行星存在的证据。他们指出：“其中一些系外行星（相对而言）小到只有月球大小，而另一些则巨大如木星。与地球不同，大多数系外行星并未紧密束缚在恒星周围，而是在太空中游荡，或在恒星之间以较为松散的方式运行。我们可以估计，这个［遥远］星系中的行星数量超过一万亿。”$^{[47]}$
\subsubsection{早期猜想}
“我们宣告这片空间是无限的……其中存在着无数与我们自身世界同类的世界。”

—— 乔尔达诺·布鲁诺（1584）$^{[48]}$

16 世纪，意大利哲学家乔尔达诺·布鲁诺是哥白尼学说（即地球和其他行星绕太阳运行的日心说）的早期支持者之一。他提出，恒星本身与太阳相似，同样也应当伴随着行星$^{[49]}$。

18 世纪，艾萨克·牛顿在其《自然哲学的数学原理》（*Principia*）结尾的《总注》（*General Scholium*）中也提及了这种可能性。他将恒星系统与太阳系作类比，写道：“如果恒星是类似系统的中心，那么它们将按照相似的设计被构造，并同样服从于唯一的主宰。”$^{[50]}$

1938 年，D. Belorizky 论证了利用凌日光度测量法搜寻“系外木星”（exo-Jupiters）的可行性$^{[51]}$。

1952 年，在第一颗“热木星”被发现之前 40 多年，奥托·斯特鲁维（Otto Struve）指出，并不存在任何必然理由认为行星不能比太阳系中的情况更靠近其母恒星。他还提出，可以通过多普勒光谱法和凌日法来探测处于短周期轨道上的超级木星行星$^{[52]}$。b6
\subsubsection{被否定的发现主张}
自 19 世纪以来，关于系外行星探测的主张便屡有提出。其中最早的一些涉及双星系统 70 蛇夫座。1855 年，东印度公司马德拉斯天文台的威廉·斯蒂芬·雅各布报告称，其轨道异常“高度可能”表明该系统中存在一个“行星体”$^{[53]}$。19 世纪 90 年代，芝加哥大学和美国海军天文台的托马斯·J. J.·西（Thomas J. J. See）声称，这些轨道异常证明在 70 蛇夫座系统中，围绕其中一颗恒星存在一个轨道周期为 36 年的暗天体$^{[54]}$。然而，福里斯特·雷·莫尔顿（Forest Ray Moulton）随后发表论文证明，在这些轨道参数下，该三体系统将高度不稳定$^{[55]}$。

关于 61 天鹅座可能拥有行星系统的主张也曾多次出现。1942 年，斯普劳尔天文台的卡伊·斯特兰德（Kaj Strand）观测到 61 天鹅座 A 与 B 的轨道运动存在微小但系统性的变化，因而推测在 61 天鹅座 A 周围必然存在一个质量约为 16 个木星质量的第三天体$^{[56]}$。此后又有多项类似主张，但较新的观测至今仍未给出确认。详见：61 Cygni：Claims of a planetary system。在研究 61 天鹅座的同时，天文学家也对拉朗德21185提出了类似的系外行星存在主张，见：Lalande 21185#Past claims of planets。

20 世纪 50 至 60 年代，斯沃斯莫尔学院的彼得·范·德·坎普（Peter van de Kamp）又发表了一系列颇具影响力的探测主张，认为巴纳德星周围存在行星$^{[57]}$。如今，天文学界普遍认为这些早期探测报告均为误判$^{[58]}$。

1991 年，安德鲁·莱恩（Andrew Lyne）、M. Bailes 和 S. L. Shemar 利用脉冲星计时变化，声称在PSR 1829−10周围发现了一颗脉冲星行星$^{[59]}$。这一消息一度引起广泛关注，但莱恩及其团队随后很快撤回了这一结论$^{[60]}$。
\subsubsection{已确认的发现}
\begin{figure}[ht]
\centering
\includegraphics[width=6cm]{./figures/7817ae950047ded9.png}
\caption{AB Pictoris 的日冕仪图像，显示其一颗伴星（左下），该天体可能是一颗褐矮星，也可能是一颗大质量行星。相关数据于 2003 年 3 月 16 日获取，使用的是甚大望远镜（VLT）上的 NACO 仪器，并在 AB Pictoris 前方加装了一个 1.4 角秒的遮挡掩膜。} \label{fig_Exopla_6}
\end{figure}
\begin{figure}[ht]
\centering
\includegraphics[width=6cm]{./figures/6c9dc65a58fd0734.png}
\caption{由海尔望远镜成像的恒星HR 8799的三颗已知行星。中央恒星的光被矢量涡旋日冕仪遮挡以消除干扰。} \label{fig_Exopla_7}
\end{figure}
\begin{figure}[ht]
\centering
\includegraphics[width=6cm]{./figures/eab4d4cff7fcea64.png}
\caption{2MASS J044144是一颗褐矮星，其伴星的质量约为木星的 5–10 倍。目前尚不清楚该伴随天体应被归类为次褐矮星还是行星。} \label{fig_Exopla_8}
\end{figure}
截至2025 年 10 月 30 日，NASA 系外行星档案库（NASA Exoplanet Archive）共列出了 6,053 颗已确认的系外行星，其中包括少数对 20 世纪 80 年代末一些存在争议的发现所作出的后续确认结果$^{[61]}$。首个在发表后获得最终确认的发现，出现在1988 年，由加拿大天文学家Bruce Campbell、G. A. H. Walker以及Stephenson Yang（分别来自维多利亚大学与不列颠哥伦比亚大学）完成$^{[62]}$。尽管他们在宣称行星探测时持相当谨慎的态度，但其径向速度观测结果表明，在恒星伽马仙王座（Gamma Cephei）周围存在一颗行星。由于当时的观测已接近仪器能力的极限，天文学界在随后数年中始终对这一发现以及其他类似观测保持怀疑态度。当时有人认为，这些“疑似行星”可能实际上是褐矮星——一种质量介于行星与恒星之间的天体。1990 年发表的补充观测支持了伽马仙王座行星存在的结论$^{[63]}$，但 1992 年的后续研究再次引发了严重质疑$^{[64]}$。最终，在2003 年，随着观测技术的改进，这颗行星的存在才得以确认$^{[65]}$。

1992 年 1 月 9 日，射电天文学家Aleksander Wolszczan与Dale Frail基于脉冲到达时间变化，宣布发现了两颗围绕毫秒脉冲星PSR 1257+12 运行的行星$^{[45]}$。这一发现随后得到确认，并通常被认为是首个无可争议的系外行星探测。后续观测进一步巩固了这一结果，而1994 年第三颗行星的确认，也使这一话题重新进入大众媒体视野$^{[66]}$。这些脉冲星行星被认为可能形成于产生该脉冲星的超新星爆发残余物之中，属于“第二轮”行星形成；或者，它们可能是某些气态巨行星在超新星爆发中幸存下来的岩石内核，随后演化到当前轨道。由于脉冲星环境极端而“暴烈”，当时人们普遍认为行星在其轨道上形成的可能性极低$^{[67]}$。

20 世纪 90 年代初，由Donald Backer领导的一组天文学家在研究一个被认为是双星脉冲星系统的天体（PSR B1620−26 b）时发现，仅靠两个天体无法解释观测到的多普勒频移，因此推断系统中必然存在第三个天体。数年内，研究人员测量了该天体对脉冲星与白矮星轨道产生的引力效应，并据此估算出其质量明显小于恒星质量。1993 年，Stephen Thorsett及其合作者正式宣布，该第三天体应为一颗行星$^{[68]}$。

1995 年 10 月 6 日，日内瓦大学的Michel Mayor与Didier Queloz宣布，在邻近的 G 型主序星飞马座 51（51 Pegasi）周围发现了一颗系外行星，这是首次明确确认一颗围绕主序星运行的系外行星$^{[69][70][71]}$。这一发现是在法国上普罗旺斯天文台完成的，被视为现代系外行星研究的开端，并最终促成二人分享2019 年诺贝尔物理学奖。随后，尤其是高分辨率光谱技术的发展，使得大量新系外行星被迅速发现——天文学家可以通过测量行星对母恒星运动产生的引力扰动来间接探测它们。后来，又通过观测行星凌星时恒星视亮度的变化，发现了更多系外行星$^{[69]}$。

最初已知的大多数系外行星，都是质量巨大、且非常靠近母恒星运行的行星。天文学家对这些所谓的“热木星”感到意外，因为传统行星形成理论认为，气态巨行星只能在远离恒星的区域形成。不过，随着观测的推进，更多类型的行星被发现，人们逐渐认识到热木星实际上只占系外行星中的少数$^{[69]}$。1999 年，仙女座 υ 星（Upsilon Andromedae）成为首个被确认拥有多颗行星的主序星系统$^{[72]}$；而Kepler-16系统则包含了首颗被发现围绕双主序星系统运行的行星$^{[73]}$。

2014 年 2 月 26 日，NASA 宣布，开普勒空间望远镜在305颗恒星周围发现并验证了 715 颗新的系外行星。这些行星是通过一种称为“多重性验证”的统计方法确认的$^{[74][75][76]}$。在此之前，大多数已确认的系外行星都是体积接近或大于木星的气态巨行星，因为它们更容易被探测；而开普勒所发现的行星，其大小多介于海王星与地球之间$^{[74]}$。

2015 年 7 月 23 日，NASA 宣布发现Kepler-452b，这是一颗尺寸接近地球、位于一颗 G2 型恒星宜居带内的行星$^{[77]}$。

2018 年 9 月 6 日，NASA 又发现了一颗距离地球约145 光年、位于室女座的系外行星$^{[78]}$。该行星被命名为Wolf 503b，其体积约为地球的两倍，围绕一颗被称为“橙矮星”的恒星运行。由于轨道极为接近母恒星，Wolf 503b 的公转周期仅约6 天。它是目前已知、位于所谓“小行星半径间隙”附近、且体积如此之大的唯一系外行星。该间隙又称为Fulton 间隙$^{[78][79]}$，指的是在观测上，很少发现半径介于地球1.5–2 倍之间的系外行星$^{[80]}$。

2020 年 1 月，科学家宣布发现TOI 700 d，这是TESS任务首次在宜居带内发现的、地球大小的行星$^{[81]}$。
\subsubsection{候选发现}
截至2020 年 1 月，NASA 的开普勒与TESS任务共识别出4,374颗尚未确认的行星候选体$^{[82]}$，其中有相当一部分尺寸接近地球、位于宜居带内，且部分环绕类太阳恒星运行$^{[83][84][85]}$。
\begin{figure}[ht]
\centering
\includegraphics[width=14.25cm]{./figures/7b0727cb5164c76f.png}
\caption{} \label{fig_Exopla_9}
\end{figure}
2020年9月，天文学家首次报告了证据，发现了一颗位于银河系外的行星——M51-ULS-1b。这颗行星通过遮挡一个明亮的X射线源（XRS）而被探测到，该X射线源位于漩涡星系（M51a）中。
\subsection{探测方法}
\begin{figure}[ht]
\centering
\includegraphics[width=14.25cm]{./figures/8f6efdc54a1400d3.png}
\caption{} \label{fig_Exopla_10}
\end{figure}
\subsubsection{直接成像}
\begin{figure}[ht]
\centering
\includegraphics[width=6cm]{./figures/6d50d76acefcc6f4.png}
\caption{从地球视角看具有近乎边缘朝向轨道的直接成像行星——$\beta$绘架座$b$} \label{fig_Exopla_11}
\end{figure}
与其母恒星相比，行星极其微弱。例如，一颗类太阳恒星的亮度大约是其所环绕的任一系外行星所反射光的十亿倍。如此微弱的光源极难被探测，而且母恒星产生的强烈眩光往往会将其淹没。为减少眩光、同时仍使行星的光可被探测，必须遮挡母恒星的光；而实现这一点是一项重大的技术挑战，需要极高水平的光学—热学稳定性 $^{[90]}$。目前所有被直接成像的系外行星，均同时具有质量很大（质量超过木星）且与其母恒星相距较远这两个特征。

专门设计的直接成像仪器，如 Gemini Planet Imager、VLT-SPHERE 和 SCExAO，将能够成像数十颗气态巨行星，但绝大多数已知的系外行星仍然只能通过间接方法被探测到。
\subsubsection{间接方法}
\textbf{凌日法}
\begin{figure}[ht]
\centering
\includegraphics[width=8cm]{./figures/1de2b37ad818bb0c.png}
\caption{当恒星位于行星之后时，其亮度看起来会减弱。} \label{fig_Exopla_12}
\end{figure}
如果一颗行星从其母恒星的盘面前方掠过（即发生凌日），那么观测到的恒星亮度就会出现微小的下降。恒星变暗的幅度取决于多种因素，包括恒星本身的大小以及行星的大小。由于凌日法要求行星的轨道必须恰好与恒星—地球的视线方向相交，因此，对于轨道取向随机的系外行星而言，被观测到发生凌日的概率相对较低。开普勒望远镜正是采用了这一方法。

\textbf{径向速度法（多普勒法）}
当行星绕恒星运行时，恒星也会绕着行星—恒星系统的质心作小幅度运动。恒星径向速度的变化——即它朝向或远离地球运动的速度——可以通过多普勒效应引起的恒星光谱线位移来探测。极其微小的径向速度变化也能够被测量到，其量级可达 $1,\mathrm{m/s}$，甚至略低于这一数值$^{[91]}$。

\textbf{凌日时刻变化}

当系统中存在多颗行星时，它们会彼此轻微扰动对方的轨道。因此，某一行星凌日时刻的细微变化，可能暗示着另一颗行星的存在，而这颗行星本身可能并不发生凌日。例如，对行星 Kepler-19b 凌日时间的变化分析表明，该系统中还存在第二颗行星——不发生凌日的 Kepler-19c$^{[92][93]}$。

\textbf{凌日持续时间变化}

当一颗行星绕多颗恒星运行，或者该行星本身拥有卫星时，其每一次凌日的发生时间都可能出现显著变化。尽管迄今为止尚未通过这一方法发现新的行星或卫星，但它已被成功用于确认许多发生凌日的环双星行星（circumbinary planets）$^{[94]}$。

\textbf{引力微透镜}
当一颗恒星的引力场像透镜一样作用时，会放大其背后遥远背景恒星的光，这一现象称为微引力透镜。若有行星绕这颗“透镜恒星”运行，它们会在放大率随时间变化的过程中造成可探测的异常。与大多数其他探测方法（通常偏向于发现近轨道行星，或在直接成像情况下偏向于远轨道的大行星）不同，微引力透镜法对探测距离类太阳恒星约 $1$–$10,\mathrm{AU}$ 处的行星最为敏感。

\textbf{天体测量法}
\begin{figure}[ht]
\centering
\includegraphics[width=8cm]{./figures/0bcd8439b7fcf6f5.png}
\caption{行星能够通过引力作用牵引其宿主恒星。} \label{fig_Exopla_13}
\end{figure}
天体测量法是通过精确测量恒星在天空中的位置，并观测其位置随时间发生的变化来实现的。由于行星的引力作用，恒星会产生可观测到的微小运动。然而，这种运动幅度极其微弱，因此直到 2020 年代，这种方法的产出一直并不高。它只带来了少数已确认的发现$^{[95][96]}$，但已被成功用于研究通过其他方法发现的行星性质。

\textbf{脉冲星计时法}

脉冲星是恒星发生超新星爆炸后留下的致密残骸，在自转过程中会以极高的稳定性发射无线电脉冲。如果有行星绕脉冲星运行，脉冲星围绕系统质心的运动会导致其与地球之间的距离随时间发生变化。结果是，来自脉冲星的无线电脉冲到达地球的时间会出现提前或延迟。这种由于脉冲星在物理上更靠近或更远离地球而产生的光行时延，称为 **罗默时延**$^{[97]}$。第一颗被确认的系外行星正是通过这种方法发现的。不过，截至 2011 年，这一方法的发现效率并不高，仅在三颗不同的脉冲星系统中探测到了五颗行星。

\textbf{变星计时法（脉动频率）}

类似于脉冲星，还有一些恒星表现出周期性的活动。其周期性的微小偏离有时可能由绕其运行的行星引起。截至 2013 年，已有少数行星通过这种方法被发现$^{[98]}$。

\textbf{反射 / 辐射调制}

当行星非常接近恒星运行时，会截获大量恒星光。随着行星绕恒星公转，由于从地球视角看到的行星相位变化，或者由于行星一侧比另一侧更炽热而产生的温度差异，接收到的总光量会发生变化$^{[99]}$。

\textbf{相对论束射效应}

相对论束射效应通过测量恒星因自身运动而产生的观测通量变化来实现。当行星使恒星更靠近或更远离地球时，恒星的亮度会发生相应变化$^{[100]}$。

\textbf{椭球形变效应}

质量较大的近轨道行星可以轻微地改变宿主恒星的形状。这会导致恒星亮度随其相对于地球的取向发生细微变化$^{[101]}$。

\textbf{偏振测量法}

在偏振测量法中，通过将行星反射的偏振光与恒星发出的非偏振光分离来进行探测。尽管尚未利用该方法发现新的行星，但已有部分已知行星通过这一方法被成功探测到$^{[102][103]}$。

\textbf{环恒星盘}

许多恒星周围存在由小行星和彗星碰撞产生的空间尘埃盘。这些尘埃会吸收恒星光并以红外辐射的形式重新释放，从而被观测到。盘中的某些结构特征可能暗示行星的存在，但这种方法通常不被视为一种确定性的行星探测手段。
\subsection{形成与演化}
行星通常会在其宿主恒星形成后的几百万年至数千万年（甚至更长）内形成$^{[104][105]}$。太阳系行星只能以其当前状态被观测到，但通过对不同年龄行星系统的观测，我们可以看到行星处在不同演化阶段的情形。现有观测对象既包括行星仍在形成中的年轻原行星盘$^{[106]}$，也包括年龄超过 100 亿年的成熟行星系统$^{[107]}$。

当行星在富含气体的原行星盘中形成时$^{[108]}$，它们会吸积由氢和氦组成的包层$^{[109][110]}$。这些包层会随时间逐渐冷却并收缩，并且根据行星质量的不同，部分或全部氢/氦最终会逃逸到太空中$^{[108]}$。这意味着，即便是类地行星，只要形成得足够早，也可能在初期具有很大的半径$^{[111][112][113]}$。

一个典型例子是Kepler-51b：它的质量仅约为地球的两倍，但半径几乎与土星相当，而土星的质量约为地球的 100 倍。Kepler-51b 的年龄只有数亿年，仍然非常年轻$^{[114]}$。
\subsection{拥有行星的恒星}
\begin{figure}[ht]
\centering
\includegraphics[width=8cm]{./figures/b627caad985cc861.png}
\caption{摩根–基南光谱分类体系} \label{fig_Exopla_14}
\end{figure}
\begin{figure}[ht]
\centering
\includegraphics[width=8cm]{./figures/4d0f6830da4a4ae8.png}
\caption{围绕双星运行的系外行星的艺术想象图 $^{[115]}$} \label{fig_Exopla_15}
\end{figure}
平均而言，每一颗恒星至少拥有一颗行星 $^{[116]}$。大约每 5 颗类太阳恒星中就有 1 颗在其宜居带内拥有一颗“类地尺寸”的行星 $^{[a][b][117]}$。

目前已知的大多数系外行星都围绕与太阳大致相似的恒星运行，即光谱型为 F、G 或 K 型的主序星。低质量恒星（红矮星，光谱型 M）较少拥有质量足够大、能够通过径向速度法探测到的行星 $^{[118][119]}$。尽管如此，借助使用凌日法、能够探测较小行星的开普勒空间望远镜，科学家已经在红矮星周围发现了数十颗行星 $^{[120]}$。

利用开普勒的数据，人们发现恒星的金属丰度与其拥有类木星大小巨行星的概率之间存在相关性：金属丰度较高的恒星比金属丰度较低的恒星更有可能拥有行星，尤其是巨行星 $^{[121]}$。

有些行星只围绕双星系统中的一颗恒星运行 $^{[122]}$，同时也已经发现了若干颗环绕双星系统中两颗恒星共同运行的环双星行星。三合星系统中已知存在少数行星 $^{[123]}$，而在四合星系统 Kepler-64 中也发现了一颗行星。
\subsection{轨道与物理参数}
主条目：系外行星的轨道与物理参数
\subsection{一般特征}
\subsubsection{颜色与亮度}
一颗行星的视亮度（视星等）取决于多个因素：观测者与行星之间的距离、行星的反照率（即反射能力），以及行星从其母恒星接收到的光量——而这又取决于行星与恒星之间的距离以及恒星本身的亮度。因此，一颗反照率较低但距离恒星很近的行星，可能看起来比一颗反照率较高但距离恒星较远的行星更亮 $^{[124]}$。
\begin{figure}[ht]
\centering
\includegraphics[width=8cm]{./figures/c4243d0399d588a6.png}
\caption{这张颜色–颜色图将太阳系行星的颜色与系外行星 HD 189733b 的颜色进行了比较。该系外行星呈现出的深蓝色是由其大气中的硅酸盐微滴所致，这些微滴会散射蓝色光。} \label{fig_Exopla_16}
\end{figure}
2013 年，人类首次测定了一颗系外行星的颜色。对 HD 189733b 的最佳拟合反照率测量表明，它呈现出深暗蓝色 $^{[125][126]}$。同年稍晚，又测定了其他几颗系外行星的颜色，其中包括在视觉上呈洋红色的 GJ 504 b $^{[127]}$，以及如果近距离观察会显得偏红色的 Kappa Andromedae b $^{[128]}$。氦行星预计在外观上呈现为白色或灰色 $^{[129]}$。

在几何反照率意义上，目前已知最暗的行星是 TrES-2b，这是一颗“热木星”，其反射的恒星光不足 1\%，反射率甚至低于煤或黑色亚克力涂料。一般认为，热木星由于其大气中含有钠和钾而应当相当昏暗，但 TrES-2b 为何如此之暗仍不清楚——这可能与某种未知的化学化合物有关 $^{[130][131][132]}$。

对于气态巨行星而言，在不存在云层显著调制的情况下，几何反照率通常会随着金属丰度或大气温度的升高而降低。云柱深度的增加会提高光学波段的反照率，但会降低某些红外波段的反照率。随着行星年龄的增加，由于云柱深度变大，其光学反照率也会随之升高。相反，随着行星质量的增加，光学反照率会降低，因为更高质量的气态巨行星具有更大的表面重力，从而导致更小的云柱深度。此外，椭圆轨道可能引起大气成分的剧烈变化，这也会产生显著影响 $^{[133]}$。

对于质量较大和／或较年轻的气态巨行星，在某些近红外波段，热辐射所占比例会超过反射光。因此，尽管光学亮度完全依赖于相位变化，但在近红外波段并不总是如此 $^{[133]}$。

气态巨行星的温度会随着时间推移以及与其母星距离的增加而降低。温度降低即使在没有云层的情况下，也会提高光学反照率。当温度降至足够低时，会形成水云，从而进一步提高光学反照率；在更低的温度下，还会形成氨云，使行星在大多数光学和近红外波段上达到最高的反照率 $^{[133]}$。
\subsubsection{磁场}
2014 年，人们通过观测氢从 HD 209458 b 行星上蒸发的方式，推断出其周围存在磁场。这是人类首次（以间接方式）探测到系外行星的磁场。该磁场的强度估计约为木星磁场的十分之一 $^{[134][135]}$。

理论上，系外行星的磁场可以通过其极光产生的射电辐射来探测，借助诸如 LOFAR 这样的高灵敏度低频射电望远镜即可实现，不过迄今尚未成功探测到此类信号 $^{[136][137]}$。这些射电辐射有望用于测量系外行星内部的自转速率，并可能提供一种比通过云层运动来推断行星自转更精确的方法 $^{[138]}$。然而，到目前为止，利用阿雷西博射电望远镜对 9 颗系外行星进行的最灵敏极光射电辐射搜索仍未取得任何发现 $^{[139]}$。

地球的磁场源于其流动的液态金属内核，但对于高压条件下的大质量“超级地球”，可能会形成不同于地球环境下的化合物。这些化合物可能具有更高的黏度和更高的熔点，从而阻止行星内部发生分层，形成没有明显内核分化的地幔结构。在超级地球所处的高压高温条件下，某些形式的氧化镁（如 MgSi(*3)O(*{12})）可能表现为液态金属，并在地幔中产生磁场 $^{[140][141]}$。

人们观测到，热木星的半径往往大于理论预期。这可能是由于恒星风与行星磁层之间的相互作用，在行星内部形成电流并产生焦耳加热，使行星受热膨胀。恒星磁活动越强，恒星风越剧烈，所产生的电流也越大，从而导致行星更强的加热和更大的膨胀。这一理论与观测到的“恒星活动性与行星半径膨胀相关”的现象相一致 $^{[142]}$。

2018 年 8 月，科学家宣布成功实现了气态氘向液态金属氢形态的转变。这一成果有助于更深入理解木星、土星及相关系外行星等气态巨行星，因为这些行星被认为内部含有大量液态金属氢，而这可能正是其强大磁场的来源 $^{[143][144]}$。

尽管此前有研究声称，近轨系外行星的磁场可能会增强其母星的耀斑活动和星斑，但 2019 年对 HD 189733 系统的研究表明这一说法并不成立。在这一研究充分的系统中未能探测到所谓的“星—行星相互作用”，从而使其他类似结论受到质疑 $^{[145]}$。随后，阿雷西博射电望远镜对 8 颗轨道半径小于 0.1 天文单位的系外行星进行射电搜索，同样未发现磁性交互作用的迹象 $^{[146]}$。

2019 年，科学家估算了 4 颗热木星的表面磁场强度，其范围约为 20–120 高斯，而木星的表面磁场强度约为 4.3 高斯 $^{[147][148]}$。

2023 年，天文学家在 YZ Ceti 行星系统中探测到可能是岩质系外行星磁层存在的首个迹象 $^{[149]}$。
\subsubsection{板块构造}
2007 年，两个独立的研究团队就较大质量超级地球是否可能存在板块构造得出了相反结论 $^{[150][151]}$：其中一方认为板块构造可能是间歇性的或停滞的 $^{[152]}$，而另一方则认为即便在行星较为干燥的情况下，超级地球上仍然极有可能发生板块构造 $^{[153]}$。

如果超级地球所含水量超过地球的 80 倍，其表面将完全被海洋覆盖，成为所谓的“海洋行星”。但若水量低于这一阈值，深部水循环仍可能在海洋与地幔之间转移足够的水分，使大陆得以存在 $^{[154][155]}$。
\subsubsection{火山活动}
55 Cancri e 表面观测到的巨大温度变化被认为可能源于火山活动释放出的大量尘埃云，这些尘埃覆盖行星表面并阻挡热辐射 $^{[156][157]}$。

由邻近行星引力牵引引起的潮汐加热，也可能促使类地系外行星产生火山活动 $^{[158][159]}$。
\subsubsection{行星环}
2007 年，恒星 V1400 Centauri 被一个天体（可能是行星或棕矮星）掩食，该天体周围被认为存在一个规模巨大的碎屑盘。该天体被命名为 J1407b，并长期被认为拥有一个比土星环还要巨大的行星环系统 $^{[160][161]}$。后续观测表明，这一结构也可能实际上是一个环行星盘 $^{[162][163]}$。

对于 HIP 41378 f，人们发现其测得的半径相对于质量而言过大，这一异常很可能是由其周围存在行星环系统所致，因此被认为有力支持其拥有行星环 $^{[164][165]}$。

在太阳系中，气态巨行星的行星环通常与行星赤道对齐。然而，对于近轨系外行星，来自母星的潮汐力会使行星最外层的环逐渐与行星绕恒星运行的轨道平面对齐，而内侧行星环仍与行星赤道对齐。如果行星自转轴存在倾角，那么内外行星环之间的不同取向将导致一个发生扭曲的环系统 $^{[166]}$。
\subsubsection{卫星}
有证据表明，其他行星周围可能存在天然卫星，通常称为系外卫星，但迄今尚无被确认的实例 $^{[citation\ needed]}$。

2012 年，人们通过行星光变曲线中的周期性变化，在 WASP-12b 周围发现了一颗候选系外卫星 $^{[167]}$，但后续观测显示，该天体实际上可能是一颗特洛伊行星 $^{[168]}$。

2013 年 12 月，在微引力透镜事件 MOA-2011-BLG-262 中发现了一颗候选系外卫星，该天体最初被认为可能是一颗绕木星质量自由漂浮行星运行、质量约为 0.5 个地球质量的系外卫星，或是一颗绕红矮星运行的海王星质量行星 $^{[169]}$。后续观测确认了后者的情形 $^{[170]}$。

2018 年 10 月 3 日，有研究报告称发现了围绕 Kepler-1625b 运行的一颗大型系外卫星的证据 $^{[171]}$；2021 年，又报告了 Kepler-1708b 周围可能存在系外卫星的证据 $^{[172]}$。不过，这些天体的真实性仍存在争议 $^{[173]}$，需要进一步观测加以确认 $^{[174]}$。

在 WASP-76b、HD 189733 b 和 WASP-49b 等热木星大气中探测到的钠元素，可能源于围绕这些行星运行的、类似木卫一（Io）的系外卫星 $^{[175][176]}$。
\subsubsection{大气层}
\begin{figure}[ht]
\centering
\includegraphics[width=8cm]{./figures/67687bcb5c6f852f.png}
\caption{两颗系外行星上晴朗与多云大气的对比 $^{[177]}$} \label{fig_Exopla_17}
\end{figure}
在若干系外行星周围已经探测到大气层。最早被观测到的例子是 HD 209458 b，于 2001 年发现。$^{[178]}$
\begin{figure}[ht]
\centering
\includegraphics[width=8cm]{./figures/528d2a470f7af9fc.png}
\caption{卡西尼号在土星卫星泰坦上进行的日落观测有助于理解系外行星的大气（艺术概念图）。} \label{fig_Exopla_19}
\end{figure}
截至 2014 年 2 月，已有超过 50 颗发生凌星的系外行星以及 5 颗被直接成像的系外行星的大气被观测到，$^{[179]}$ 由此实现了分子光谱特征的探测、昼夜温度梯度的观测，以及对大气垂直结构的约束。$^{[180]}$ 此外，在一颗未发生凌星的“热木星”——金牛座 τ（Tau Boötis）b 上，也探测到了大气层。$^{[181][182]}$

2017 年 5 月，科学家发现，从距离地球约一百万英里的轨道卫星上看到的、仿佛闪烁的地球反射光，实际上来自大气中冰晶的反射。$^{[183][184]}$ 用于判定这一现象的技术，可能有助于研究遥远世界（包括系外行星）的大气。$^{[citation\ needed]}$

\textbf{彗星状尾迹}

开普勒-1520b 是一颗体积较小的岩质行星，距离其恒星极近，正在发生蒸发，并在其后方留下类似彗星的云尘尾迹。$^{[185]}$ 这些尘埃可能是火山喷发产生的火山灰，由于该行星表面重力较小而逃逸；也可能来源于金属在高温下被汽化，随后金属蒸气冷凝成尘埃。$^{[186]}$

2015 年 6 月，科学家报告称，GJ 436 b 的大气正在蒸发，在行星周围形成了一个巨大的气体云，并在宿主恒星辐射作用下，产生了一条长达 1400 万千米（900 万英里）的拖曳尾迹。$^{[187]}$
\subsubsection{日照分布模式}
处于 1:1 自转—公转共振（潮汐锁定）的行星，其恒星将始终直射行星表面的同一位置，该区域极为炽热，而对侧半球则完全得不到光照、极端寒冷。这样的行星外观可能类似“眼球”，其中热点相当于“瞳孔”。$^{[188]}$ 具有偏心轨道的行星可能会锁定在其他共振状态中，例如 3:2 或 5:2 共振，这将产生“双眼球”式的日照模式，在东西两个半球各有一个热点。$^{[189]}$ 若行星同时具有偏心轨道和倾斜的自转轴，其日照分布模式将更加复杂。$^{[190]}$
\subsection{表面}
\subsubsection{表面成分}
通过将发射与反射光谱同透射光谱进行比较，可以区分表面特征与大气特征。系外行星的中红外光谱可能探测到岩质表面，而近红外光谱则可能识别岩浆海或高温熔岩、含水硅酸盐表面以及水冰，从而提供一种明确区分岩质行星与气态行星的方法。$^{[191]}$
\subsubsection{表面温度}
\begin{figure}[ht]
\centering
\includegraphics[width=8cm]{./figures/e50cf522d09216a7.png}
\caption{系外行星大气中温度逆温现象的艺术示意图。$^{[192]}$} \label{fig_Exopla_20}
\end{figure}
通过测量一颗行星从其母恒星接收到的光辐射强度，可以估算其温度。例如，行星 OGLE-2005-BLG-390Lb 的表面温度被估计约为 −220 °C（50 K）。然而，这类估计可能存在相当大的误差，因为它们依赖于行星通常未知的反照率，而且诸如温室效应等因素也可能引入额外的不确定性。一些行星的温度已经通过观测其在轨道运动过程中红外辐射的变化、以及被母恒星掩食时的辐射差异而得到测量。例如，行星 HD 189733b 的昼半球平均温度估计为 1205 K（932 °C），而夜半球约为 973 K（700 °C）。$^{[193]}$
\subsection{宜居性}
随着越来越多行星被发现，系外行星学这一领域正不断发展，对太阳系外世界的研究也日益深入，并最终将触及太阳系之外行星上是否存在生命这一前景问题。$^{[194]}$ 在宇宙尺度的距离上，只有当生命在行星尺度上充分发展，并对行星环境产生了强烈改造，且这些改造无法用经典物理化学过程（即非平衡过程）来解释时，生命才可能被探测到。$^{[194]}$ 例如，地球大气中的分子氧（O$_2$）是植物和多种微生物进行光合作用的结果，因此可作为系外行星上生命存在的一个指示；尽管如此，少量氧气也可能通过非生物过程产生。$^{[195]}$ 此外，一颗潜在宜居的行星必须绕一颗稳定的恒星运行，并位于这样一个距离范围内：在该范围中，具有行星质量的天体在足够的大气压条件下能够在其表面维持液态水的存在。$^{[196][197]}$
\subsubsection{宜居带}
\begin{figure}[ht]
\centering
\includegraphics[width=8cm]{./figures/f26e2c5965b952b2.png}
\caption{一幅示意图，展示了基于 2024 年 9 月数据，不同恒星类型下宜居带边界的分布情况。图中同时标出了地球以及位于宜居带内的 42 颗潜在类岩质系外行星。} \label{fig_Exopla_21}
\end{figure}
恒星周围的**宜居带**是指温度恰到好处、能够使行星表面存在液态水的区域；也就是说，既不能离恒星太近以致水分蒸发，也不能离得太远而使水分冻结。恒星产生的热量会随其大小和年龄而变化，因此不同恒星的宜居带距离也各不相同。此外，行星自身的大气条件会影响其保持热量的能力，因此宜居带的位置还取决于行星的类型：沙漠行星（又称干燥行星）由于水量极少，其大气中的水汽含量低于地球，温室效应较弱，这意味着沙漠行星可以在比地球—太阳距离更靠近其恒星的位置维持局部的液态水绿洲。水量不足也意味着反射热量回太空的冰更少，因此沙漠行星宜居带的外边界会更远$^{[198][199]}$。具有厚重氢气大气层的岩石行星，甚至可以在远远超过地球—太阳距离的位置维持表面液态水$^{[200]}$。质量更大的行星往往拥有更宽的宜居带，因为更强的引力会降低水云柱的厚度，从而减弱水汽的温室效应，使宜居带的内边界向恒星方向移动$^{[201]}$。

行星自转速率是决定大气环流以及云层分布格局的主要因素之一：自转较慢的行星更容易形成厚重的云层，反射更多恒星辐射，因此可以在更靠近恒星的位置保持宜居。如果地球保持当前的大气条件，但拥有金星那样缓慢的自转速度，那么即使位于金星轨道上也可能是宜居的。若金星因失控温室效应而失去了其水洋，那么它在过去很可能具有更高的自转速率；另一种可能是，金星在形成过程中水汽直接逃逸到太空，从未拥有过海洋，并且在其整个历史中始终保持缓慢自转$^{[202][203]}$。

潮汐锁定行星（亦称“眼球行星”$^{[204]}$）由于云层效应的存在，其宜居距离可能比此前认为的更靠近恒星：在高恒星辐照条件下，强烈对流会在亚恒星点附近形成厚重的水云，大幅提高行星反照率并降低地表温度$^{[205]}$。

位于低金属丰度恒星宜居带内的行星，相较于高金属丰度恒星，更有利于陆地复杂生命的出现。这是因为高金属丰度恒星的光谱不利于臭氧的形成，从而使更多紫外辐射到达行星表面$^{[206][207]}$。

传统上，宜居带通常是依据行星表面温度来定义的；然而，地球上超过一半的生物量来自地下微生物$^{[208]}$，而温度会随深度增加。因此，即便行星表面冻结，地下环境仍可能适宜微生物生存。若将这一点纳入考虑，宜居带可以延伸到距离恒星更远的位置$^{[209]}$，甚至流浪行星在足够深的地下也可能存在液态水$^{[210]}$。在宇宙更早的时期，宇宙微波背景辐射的温度甚至可能使任何存在的岩石行星在不论其与恒星距离多远的情况下，都能够在表面维持液态水$^{[211]}$。类似木星的气态巨行星本身或许并不宜居，但它们可能拥有宜居的卫星$^{[212]}$。
\subsubsection{冰期与雪球态}
宜居带的外缘是行星会被完全冻结的位置，但即使位于宜居带内部较深处的行星，也可能周期性进入冻结状态。如果轨道变化或其他因素导致降温，就会形成更多冰层；而冰会反射阳光，进一步加剧降温，从而形成正反馈回路，直到行星完全或几乎完全冻结。当行星表面被冻结时，二氧化碳的风化作用停止，火山释放的二氧化碳便会在大气中累积，增强温室效应，最终再次解冻行星。
具有较大轴倾角的行星$^{[213]}$不太容易进入雪球态，并且能够在距离恒星更远的位置保持液态水。轴倾角的大幅波动甚至可能比固定的大轴倾角产生更强的增温效应$^{[214][215]}$。
颇具悖论的是，围绕较冷恒星（如红矮星）运行的行星反而不太容易进入雪球态，因为冷恒星发出的红外辐射主要集中在会被冰吸收的波段，从而加热冰层$^{[216][217]}$。
\subsubsection{潮汐加热}
如果行星具有偏心轨道，除恒星辐射之外，潮汐加热还能提供另一种能量来源。这意味着位于辐射宜居带**内的偏心行星可能会过热，反而不利于液态水的存在。随着时间推移，潮汐作用会使轨道趋于圆化，因此，宜居带中可能存在一些圆轨道但缺乏水的行星——它们在过去曾具有偏心轨道$^{[218]}$。位于宜居带之外更远处的偏心行星，其表面仍可能被冻结，但潮汐加热有可能在地下形成类似木卫二（Europa）的次表层海洋$^{[219]}$。在某些行星系统中（例如仙女座上星系统 Upsilon Andromedae），轨道偏心率会受到系统内其他行星的扰动而得以维持，甚至发生周期性变化。潮汐加热还可能引发地幔逸气，从而促进大气的形成与补充$^{[220]}$。
\subsubsection{潜在宜居行星}
2015 年的一项综述将Kepler-62f、Kepler-186f 和 Kepler-442b认定为最有可能具备宜居性的系外行星候选者$^{[221]}$。它们分别距离地球约1000、490 和 1120 光年。其中，Kepler-186f 的大小与地球相近，其半径约为地球的1.2 倍，并位于其红矮星宿主的宜居带外缘附近$^{[222]}$。在最近的类地系外行星候选者中，比邻星 b距离地球约4.2 光年，其平衡温度估计为−39 °C（234 K）$^{[223]}$。
\begin{figure}[ht]
\centering
\includegraphics[width=6cm]{./figures/b14746f83031ca05.png}
\caption{} \label{fig_Exopla_22}
\end{figure}
\textbf{地球大小行星}

\begin{itemize}
\item 2013 年 11 月的一项估计表明，银河系中约$22\pm8\%$*的类太阳恒星$^{[a]}$可能在其宜居带$^{[c]}$内拥有一颗地球大小$^{[b]}$的行星$^{[4][117]}$。若假设银河系中约有2000 亿颗恒星$^{[d]}$，则潜在宜居的“地球”数量约为110 亿；若将红矮星也计入，则可增至 400 亿$^{[6]}$。

\item Kepler-186f：2014 年 4 月公布，位于一颗红矮星的宜居带内，行星半径约为地球的1.2 倍。
\item 比邻星 b（Proxima Centauri b）：位于比邻星的宜居带内；比邻星是距离太阳系最近的已知恒星，该行星的估计最小质量约为地球的1.27 倍。
\item 2013 年 2 月，研究人员推测多达6\% 的小型红矮星可能拥有地球大小行星。这意味着距离太阳系最近的此类行星可能仅13 光年；在95\% 置信区间下，估计距离增至21 光年$^{[224]}$。2013 年 3 月的一项修订估计将红矮星宜居带内地球大小行星的出现率提高至 50\%$^{[225]}$。
\item Kepler-452b：其半径约为地球的1.63 倍，是首个发现的、位于G2 型类太阳恒星宜居带内的近地球大小行星（2015 年 7 月）$^{[226]}$。
\end{itemize}
\begin{figure}[ht]
\centering
\includegraphics[width=6cm]{./figures/8a16152038425d31.png}
\caption{开普勒任务发现的、位于其母星宜居带内的小型行星对比图} \label{fig_Exopla_23}
\end{figure}
\subsection{行星系统}
系外行星通常是围绕恒星运行的多行星系统成员。行星之间通过引力相互作用，有时会形成共振系统，使行星的轨道周期呈整数比关系。例如，开普勒-223 系统包含四颗行星，处于 8:6:4:3 的轨道共振之中$^{[227]}$。

一些“热木星”绕其母星运行的方向与母星自转方向相反$^{[228]}$。一种可能的解释是，热木星往往形成于致密星团中，在那里扰动更为频繁，行星也更有可能被邻近恒星通过引力俘获$^{[229]}$。
\subsection{搜索项目}
\begin{itemize}
\item ANDES —— ArmazoNes 高分辨率阶梯光栅光谱仪，是一台用于行星发现与行星性质刻画的光谱仪，预计将安装在欧洲南方天文台（ESO）的 39.3 米口径 ELT 望远镜上。ANDES 早先称为 HIRES，而 HIRES 本身是由此前的 CODEX（光学高分辨率）和 SIMPLE（近红外高分辨率）两种光谱仪概念的团队合并后形成的$^{[230]}$。
\item CoRoT —— 发现第一颗凌日型岩质行星的空间望远镜$^{[231]}$。
\item ESPRESSO —— 一台用于搜寻岩质行星并进行高稳定度光谱观测的光谱仪，安装在 ESO 位于智利北部阿塔卡马沙漠帕瑞纳尔山顶的甚大望远镜（VLT，4 × 8.2 米）上。
\item HARPS —— 高精度阶梯光栅行星搜寻光谱仪，安装在 ESO 位于智利拉西拉天文台的 3.6 米望远镜上$^{[232]}$。
\item Kepler —— 采用凌日法搜寻大量系外行星的空间任务$^{[120]}$。
\item TESS —— 搜索新系外行星的空间任务，运行时间为 2018 至 2020 年，并通过转向观测覆盖全天空恒星。截至 2025 年 3 月 22 日，TESS 已识别出 7,525 个系外行星候选体，其中 618 个已被确认$^{[233]}$。
\end{itemize}
\subsection{参见}
\begin{itemize}
\item 从遥远的恒星系统探测地球
\item 系外行星百科
\item 小说中的系外行星
\item 复杂生命宜居带
\item 系外行星列表
\item NASA系外行星档案
\item 行星捕获
\end{itemize}
\subsection{注释}\\
a.就这一1/5的统计数据而言，“类太阳”指的是G型恒星。由于缺乏类太阳恒星的数据，该统计数据是基于K型恒星.\\
b.就这一1比5的统计数据而言，地球大小指1至2个地球半径。\\
c.就这一1比5的统计数据而言，“宜居带”指恒星辐射通量为地球的0.25至4倍的区域（对应于太阳的0.5至2天文单位）。\\
d.大约四分之一的恒星是类GK型太阳。银河系中恒星的总数尚不确切，但若以总计2000亿颗恒星来估算，银河系中大约有500亿颗类GK型太阳，其中约五分之一（22\%，即110亿颗）的恒星周围可能存在位于宜居带的类地行星。如果将红矮星也考虑在内，这一数字将增至400亿颗。
\subsection{参考文献}
\begin{enumerate}
\item “行星系统综合数据”。美国国家航空航天局系外行星档案。检索日期：2025年12月4日
\item 布伦南，帕特（2022年3月21日）。“宇宙里程碑：美国国家航空航天局确认发现5000颗系外行星”。美国国家航空航天局。检索日期：2022年4月2日
\item 巴列斯特罗斯，F. J.；费尔南德斯-索托，A.；马丁内斯，V. J.（2019）。“标题：深入探索系外行星：水海是最常见的吗？”天体生物学. 19 (5): 642–654. doi:10.1089/ast.2017.1720. hdl:10261/213115. PMID 30789285. S2CID 73498809.
\item 桑德斯，R.（2013年11月4日）。“天文学家解答关键问题：宜居行星有多常见？”. newscenter.berkeley.edu.
\item 佩蒂古拉，E. A.；霍华德，A. W.；马西，G. W.（2013）。“围绕类太阳恒星运行的地球大小行星的普遍性”. 美国国家科学院院刊. 110 (48): 19273–19278. arXiv:1311.6806. 天文学文献标识码:2013PNAS..11019273P. DOI:10.1073/pnas.1319909110. PMC 3845182. PMID 24191033.
\item 汗，阿米娜（2013年11月4日）。“银河系可能孕育数十亿颗类地行星”。洛杉矶时报。检索日期：2013年11月5日
\item "HR 2562 b". Caltech. Retrieved 15 February 2018.
\item 科诺帕基，奎因·M.；拉莫，朱利安；杜谢纳，加斯帕尔；菲利帕佐，约瑟夫·C.；乔尔拉·戈德弗里，佩吉·A.；马鲁瓦，克里斯蒂安；尼尔森，埃里克·L.（2016年9月20日）。“发现邻近碎屑盘恒星HR 2562的亚恒星伴星” （PDF）. 天体物理学杂志快报. 829 (1): 10. arXiv:1608.06660. Bibcode:2016ApJ...829L...4K. doi:10.3847/2041-8205/829/1/L4. hdl:10150/621980. S2CID 44216698.
\item 迈尔，A.；罗德，L.；拉佐尼，C.；博卡莱蒂，A.；布兰德纳，W.；加利歇，R.；康塔卢贝，F.；梅萨，D.；克拉尔，H.；博斯特，H.；肖万，G.；德西德拉，S.；扬森，M.；凯普勒，M.；奥洛夫松，J.；奥热罗，J.；达姆根，S.；亨宁，T.；泰博，P.；邦努瓦，M.；费尔特，M.；格拉顿，R.；拉格朗日，A.；朗格卢瓦，M.；迈耶，M. R.；维冈，A.；多拉齐，V.；哈格尔贝格，J.；勒科罗莱，H.；利吉，R.；鲁昂，D.；萨姆兰德，M.；施密特，T.；于德里，S.；祖尔洛，A.；阿贝，L.；卡勒，M.；德尔布尔贝，A.；福蒂耶，P.；马尼亚尔，Y.；莫雷尔，D.；穆兰，T.；巴甫洛夫，A.；佩雷，D.；佩蒂，C.；拉莫斯，J. R.；里加尔，F.；鲁，A.；韦伯，L.（2018）。“VLT/SPHERE天体测量确认与棕矮星伴星HR 2562 B的轨道分析”. 天文学与天体物理学. 615：A177。arXiv:1804.04584. Bibcode:2018A&A...615A.177M. doi:10.1051/0004-6361/201732476.
\item 博登海默，彼得；丹杰洛，詹纳罗；利绍尔，杰克·J.；福特尼，乔纳森·J.；索蒙，迪迪埃（2013）。“通过核心成核吸积形成的巨型行星和低质量棕矮星中的氘燃烧”。天体物理学杂志. 770 (2): 120. arXiv:1305.0980. 文献标识码:2013ApJ...770..120B. DOI:10.1088/0004-637X/770/2/120. S2CID 118553341.
\item 安格拉达-埃斯库德，吉列姆；阿马多，佩德罗·J.；巴恩斯，约翰；等（2016）。“比邻星周围宜居轨道上的一颗类地行星候选体”. 自然. 536 (7617): 437–440. arXiv:1609.03449. Bibcode:2016Natur.536..437A. doi:10.1038/nature19106. PMID 27558064. S2CID 4451513.
\item 扎科斯，伊莱恩（2018年2月5日）。“银河系外可能存在超过一万亿颗行星”。国家地理于2021年4月28日存档自原始网页。检索日期：2018年2月5日
\item 曼德尔鲍姆，瑞安·F.（2018年2月5日）。“科学家在遥远星系中发现数千颗行星的证据”。Gizmodo。检索日期：2018年2月5日
\item 奥布赖，丹尼斯（2015年1月6日）。“随着类地行星数量增多，天文学家开始思考下一步该何去何从”. 《纽约时报》。于2022年1月1日存档自原文。
\item 奥卡拉汉，乔纳森（2023年1月23日）。“詹姆斯·韦伯太空望远镜预示着系外行星科学的新纪元——詹姆斯·韦伯太空望远镜正在为系外行星研究以及寻找地球以外的生命开启一个令人振奋的新篇章”。《科学美国人》。检索于2023年1月23日
\item 贝希曼，C.；杰利诺，克里斯托弗·R.；柯克帕特里克，J. 戴维；库欣，迈克尔·C.；多德森-罗宾逊，萨莉；马利，马克·S.；莫利，卡罗琳·V.；赖特，E. L.（2014）。“WISE Y型矮星作为褐矮星与系外行星关联的探针”。天体物理学杂志. 783 (2): 68. arXiv:1401.1194. Bibcode:2014ApJ...783...68B. doi:10.1088/0004-637X/783/2/68. S2CID 119302072.
\item 德雷克，纳迪娅（2014年3月13日）。“银河系中孤独行星指南”。国家地理。于2021年5月18日从原文存档。检索日期：2022年1月17日
\item 斯特里加里，L. E.；巴尔纳贝，M.；马歇尔，P. J.；布兰福德，R. D.（2012）。“银河系的游牧者”。皇家天文学会月报423（2）：1856–1865。arXiv：1201.2687天文学文献标识码2012MNRAS.423.1856SDOI10.1111/j.1365-2966.2012.21009.xS2CID119185094估计，在质量介于0.08至1倍太阳质量的主序星中，每颗恒星周围约有700个质量大于10倍太阳质量的天体（大致相当于火星的质量），而银河系中此类恒星的数量高达数十亿颗。
\item “国际天文学联合会2006年大会：国际天文学联合会决议投票结果”，2006年。于2024年12月25日从原文存档。检索日期：2025年3月31日。
\item 布里特，R. R.（2006）。“为什么行星永远不会被定义”。Space.com。检索日期：2008年2月13日
\item “太阳系外行星工作组：‘行星’的定义”。国际天文学联合会立场声明。2003年2月28日。于2021年4月19日从原文存档。检索日期：2025年3月31日
\item “系外行星的官方工作定义”。国际天文学联合会立场声明。于2025年2月12日从原始页面存档。检索日期：2025年3月31日
\item 勒卡韦耶·德斯昂日，A.；利绍尔，杰克·J.（2022年6月）。“国际天文学联合会关于系外行星的工作定义”. 《新天文评论》. 94 101641. arXiv:2203.09520. 文献标识码:2022NewAR..9401641L. DOI:10.1016/j.newar.2022.101641. S2CID 247065421.
\item 莫尔达西尼，C.；阿利贝尔，扬；本茨，威利；纳埃夫，多米尼克（2008）。“通过核心吸积形成巨行星”。极端太阳系. 398: 235. arXiv:0710.5667. Bibcode:2008ASPC..398..235M.
\item 巴拉夫，I.；沙布里耶，G.；巴尔曼，T.（2008）。“超地球至超木星系外行星的结构与演化。Ⅰ. 内部重元素富集”。天文学与天体物理学. 482 (1): 315–332. arXiv:0802.1810. Bibcode:2008A&A...482..315B. doi:10.1051/0004-6361:20079321. S2CID 16746688.
\item 布希，弗朗索瓦；埃布拉尔，纪尧姆；于德里，斯蒂芬；德尔福斯，哈维耶；博瓦，伊莎贝尔；德索尔，摩根；邦菲，泽维尔；埃根贝格尔，安妮；埃伦赖希，大卫；福尔韦耶，蒂埃里；勒科罗勒，埃尔韦；拉格兰日，安妮-玛丽；洛维斯，克里斯托夫；穆图，克莱尔；佩佩，弗朗切斯科；佩里耶，克里斯蒂安；蓬，弗雷德里克；凯洛兹，迪迪埃；桑托斯，努诺·C.；塞格兰桑，达米安；维达尔-马贾尔，阿尔弗雷德（2009）。“SOPHIE北天系外行星研究。I. 一颗靠近行星/棕矮星过渡区的伴星，围绕HD16760”. 天文学与天体物理学. 505 (2): 853–858. 文献标识码:2009A&A...505..853B. DOI:10.1051/0004-6361/200912427.
\item 库马尔，希夫·S.（2003）。“命名法：棕矮星、气态巨行星与？”。棕矮星. 211: 532. 文献标识码:2003IAUS..211..529B.
\item 布兰特，T. D.；麦克埃尔温，M. W.；特纳，E. L.；梅德，K.；施皮格尔，D. S.；久住原，M.；施利德，J. E.；维斯涅夫斯基，J. P.；阿部，L.；比勒，B.；布兰德纳，W.；卡森，J.；柯里，T.；埃格纳，S.；费尔特，M.；戈洛塔，T.；后藤，M.；格雷迪，C. A.；居永，O.；桥本J；早野Y；林M；林S；亨宁T；霍达普K W；犬冢S；石井M；伊江M；扬森M；神取R；等（2014）。“对行星种子及其他高对比度系外行星巡天的统计分析：大质量行星还是低质量棕矮星？”天体物理学报. 794 (2): 159. arXiv:1404.5335. Bibcode:2014ApJ...794..159B. doi:10.1088/0004-637X/794/2/159. S2CID 119304898.
\item 施劳夫曼，凯文·C.（2018年1月22日）。“行星质量上限的证据及其对巨行星形成的启示”。《天体物理学杂志》853（1）：37。arXiv：1801.06185文献标识码2018ApJ...853...37SDOI10.3847/1538-4357/aa961cISSN1538-4357S2CID55995400
\item 斯皮格尔，D. S.；伯罗斯，亚当；米尔索姆，J. A.（2011）。“褐矮星与巨行星的氘燃烧质量极限”。天体物理学报. 727 (1): 57. arXiv:1008.5150. Bibcode:2011ApJ...727...57S. doi:10.1088/0004-637X/727/1/57. S2CID 118513110.
\item 施奈德，J.；德迪厄，C.；勒西达内，P.；萨瓦勒，R.；佐洛图金，I.（2011）。“系外行星的定义与编目：exoplanet.eu数据库”。天文学与天体物理学. 532(79)：A79。arXiv:1106.0586. Bibcode:2011A&A...532A..79S. doi:10.1051/0004-6361/201116713. S2CID 55994657.
\item 施奈德，让（2016）。“第III.8章：系外行星与棕矮星——科罗特的视角及未来展望”。系外行星与棕矮星：科罗特的视角及未来展望。第157页。arXiv:1604.00917. doi:10.1051/978-2-7598-1876-1.c038. ISBN 978-2-7598-1876-1. S2CID 118434022.
\item 哈策斯·海克·劳尔，阿蒂·P.（2015）。“基于质量-密度关系的巨行星定义”。天体物理学杂志. 810(2)：L25。arXiv:1506.05097. 文献标识码:2015ApJ...810L..25H. doi:10.1088/2041-8205/810/2/L25. S2CID 119111221.
\item 赖特，J. T.；法库里，O.；马西，G. W.；韩，E.；冯，Y.；约翰逊，约翰·阿舍；霍华德，A. W.；费舍尔，D. A.；瓦伦蒂，J. A.；安德森，J.；皮斯库诺夫，N.（2010）。“系外行星轨道数据库”。太平洋天文学会出版物. 123 (902): 412–422. arXiv:1012.5676. Bibcode:2011PASP..123..412W. doi:10.1086/659427. S2CID 51769219.
\item “系外行星纳入系外行星档案的判定标准”.exoplanetarchive.ipac.caltech.edu。检索日期：2022年1月17日
\item 巴斯里，吉博；布朗，迈克尔·E.（2006）。“从原行星到棕矮星：什么是行星？” （PDF）. 地球与行星科学年鉴（已提交稿件）。34: 193–216. arXiv:astro-ph/0608417. 文献标识码:2006AREPS..34..193B. DOI:10.1146/annurev.earth.34.031405.125058. S2CID 119338327.
\item 莱伯特，詹姆斯（2003）。“命名法：棕矮星、气态巨行星与？”棕矮星. 211: 533. 文献标识码:2003IAUS..211..529B.
\item “新型深度学习方法为开普勒望远镜的行星总数新增301颗行星——美国国家航空航天局”，2021年11月23日。
\item 《系外行星目录》。1995年。
\item “ESO的SPHERE首次揭示一颗系外行星”.www.eso.org。检索日期：2017年7月7日
\item “国际天文学联合会 | IAU”.www.iau.org。检索日期：2017年1月29日
\item 法里希，J.（2016年4月1日）。“白矮星周围的尘埃碎屑与污染”。《新天文学评论》71：9–34。arXiv1604.03092文献标识码2016NewAR..71....9FDOI10.1016/j.newar.2016.03.001ISSN1387-6473
\item 兰道，伊丽莎白（2017年11月1日）。“被忽视的瑰宝：系外行星的首个证据”。美国国家航空航天局。检索日期：2017年11月1日
\item 《太阳系百科全书》，第三版，2014年，第963页，蒂尔曼·施波恩、多丽丝·布罗伊尔、托伦斯·约翰逊
\item 沃尔什赞，A.；弗雷尔，D. A.（1992）。“毫秒脉冲星PSR1257+12周围的一个行星系统”。自然. 355 (6356): 145–147. 天文学编码:1992Natur.355..145W. DOI:10.1038/355145a0. S2CID 4260368.
\item “预先生成的系外行星图谱”.exoplanetarchive.ipac.caltech.eduNASA系外行星档案. 检索日期2023年7月10日
\item 扎霍斯，埃莱娜（2018年2月5日）。“银河系外可能存在着超过一万亿颗行星——一项新研究首次提供了银河系外存在系外行星的证据”。国家地理学会。于2021年4月28日从原文存档。检索于2022年7月31日
\item 埃利·马奥尔（1987）。“第24章：新宇宙学”。迈向无穷与超越：无限的文化史。最初发表于《论无限的宇宙与世界》[De l'infinito universo et mondi]，作者乔尔达诺·布鲁诺（1584）。马萨诸塞州波士顿：比尔克豪瑟出版社。第198页ISBN978-1-4612-5396-9
\item “乔尔达诺·布鲁诺 | 传记、逝世与事实 | 百科全书”.www.britannica.com. 2025年7月5日. 检索日期：2025年8月10日
\item 牛顿，艾萨克；I. 伯纳德·科恩；安妮·惠特曼（1999）[1713]。《原理：全新译本与指南》加州大学出版社。第940页。ISBN 978-0-520-08816-0.
\item 太阳，变星, D.贝洛里茨基，1938
\item 斯特鲁维，奥托（1952）。“高精度恒星径向速度研究项目建议”。天文台. 72: 199–200. 文献标识码:1952Obs....72..199S.
\item 雅各布，W. S.（1855）。“关于双星蛇夫座70星所呈现的某些异常现象”. 皇家天文学会月报. 15 (9): 228–230. 文献标识码:1855MNRAS..15..228J. DOI:10.1093/mnras/15.9.228.
\item 见，T. J. J.(1896)。“关于70蛇夫座的轨道研究，以及由一颗不可见天体作用引起的该系统运动中的周期性摄动”。天文期刊. 16: 17–23. 文献标识码:1896AJ.....16...17S. DOI:10.1086/102368.
\item 谢里尔，T. J.（1999）。“充满争议的职业生涯：T. J. J. 希的异常现象” （PDF）. 天文学史期刊. 30 (98): 25–50. 文献标识码:1999JHA....30...25S. DOI:10.1177/002182869903000102. S2CID 117727302.
\item 斯特兰德，K. Aa. （1943）。“61天鹅座作为三重系统”. 太平洋天文学会出版物. 55 (322): 29–32. 文献标识码:1943PASP...55...29S. DOI:10.1086/125484.
\item 范德坎普，P.（1969）。“巴纳德恒星的替代动力学分析”。天文学杂志. 74: 757–759. 文献标识码:1969AJ.....74..757V. DOI:10.1086/110852.
\item 老板，艾伦（2009）。拥挤的宇宙：寻找宜居行星。基本图书。第31–32. ISBN 978-0-465-00936-7.
\item 贝尔斯，M.；莱恩，A. G.；谢马尔，S. L.（1991）。“一颗围绕中子星PSR1829–10运行的行星”。自然. 352 (6333): 311–313. 文献标识码:1991Natur.352..311B. DOI:10.1038/352311a0. S2CID 4339517.
\item 莱恩，A. G.；贝尔斯，M.（1992）。“没有行星围绕PS R1829–10运行”. 自然. 355 (6357): 213. 文献标识码:1992Natur.355..213L. DOI:10.1038/355213b0. S2CID 40526307.
\item 施奈德，J.“互动式太阳系外行星目录”.太阳系外行星百科全书. 检索日期：2025年10月30日
\item 坎贝尔，B.；沃克，G. A. H.；杨，S.（1988）。“对类太阳恒星亚恒星伴星的搜寻”. 天体物理学报. 331: 902. 文献标识码:1988ApJ...331..902C. doi:10.1086/166608.
\item 劳顿，A. T.；赖特，P.（1989）。“仙后座γ星的行星系统？”英国行星际学会期刊. 42: 335–336. 文献标识码:1989JBIS...42..335L.
\item 沃克，G. A. H；博伦德，D. A.；沃克，A. R.；欧文，A. W.；杨，S. L. S.；拉尔森，A.（1992）。“仙后座γ——是自转还是行星伴星？”天体物理学报快报396（2）：L91－L94文献标识码：1992ApJ...396L..91WDOI10.1086/186524
\item 哈策斯，A. P.；科克兰，威廉·D.；恩德尔，迈克尔；麦克阿瑟，芭芭拉；保尔森，黛安·B.；沃克，戈登·A. H.；坎贝尔，布鲁斯；杨，斯蒂芬森（2003）。“仙后座γ A的行星伴星”。天体物理学杂志. 599 (2): 1383–1394. arXiv:astro-ph/0305110. 文献标识码:2003ApJ...599.1383H. DOI:10.1086/379281. S2CID 11506537.
\item 霍尔茨，罗伯特（1994年4月22日）。“科学家发现恒星周围存在新行星的证据”。洛杉矶时报，经由科技在线。于2013年5月17日存档自原文。检索日期：2012年4月20日
\item 罗德里格斯·巴克罗，奥斯卡·奥古斯托（2017）。太阳系之外的人类存在 [Human presence beyond the solar system] （西班牙语）。RBA。第29页。ISBN 978-84-473-9090-8.
\item “已确认最古老的行星”.哈勃网站。检索日期：2006年5月7日
\item 温兹，约翰（2019年10月10日）。“来自炽热怪异行星的启示”。可知晓杂志。年度评论。DOI：10.1146/knowable-101019-2。检索日期：2022年4月4日
\item 梅耶尔，M.；凯洛兹，D.（1995）。“一颗类太阳恒星的木星质量伴星”。自然. 378 (6555): 355–359. 天文学编码:1995Natur.378..355M. DOI:10.1038/378355a0. S2CID 4339201.
\item 吉布尼，伊丽莎白（2013年12月18日）。“寻找姐妹地球”. 自然. 504 (7480): 357–365. Bibcode:2013Natur.504..357。. doi:10.1038/504357a. PMID 24352276.
\item 利索尔，J. J.（1999）。“仙女座υ星系的三颗行星”. 自然. 398 (6729): 659. 文献标识码:1999Natur.398..659L. DOI:10.1038/19409. S2CID 204992574.
\item 多伊尔，L. R.；卡特，J. A.；法布里基，D. C.；斯劳森，R. W.；霍威尔，S. B.；温恩，J. N.；奥罗兹，J. A.；普尔沙，A.；韦尔什，W. F.；奎因，S. N.；拉瑟姆，D.；托雷斯，G.；布赫哈维，L. A.；马西，G. W.；福特尼，J. J.；什波雷尔，A.；福特，E. B.；利萨uer，J.J. 拉戈齐尼，D. 拉戈齐尼，M. 鲁克，N. 巴塔利亚，J. M. 詹金斯，W. J. 博鲁基，D. 科赫，C. K. 米多尔，J. R. 霍尔，S. 麦考利夫，M. N. 法内利，E. V. 昆塔纳，M. J. 霍尔曼等（2011）。“开普勒-16：一颗凌星的环双星行星”。科学. 333 (6049): 1602–1606. arXiv:1109.3432. Bibcode:2011Sci...333.1602D. doi:10.1126/science.1210923. PMID 21921192. S2CID 206536332.
\item 约翰逊，米歇尔；哈林顿，J.D.（2014年2月26日）。“美国国家航空航天局开普勒任务宣布发现大量行星，新增715颗新世界”。美国国家航空航天局。于2014年2月26日存档自原文。检索日期：2014年2月26日
\item 沃爾，迈克（2014年2月26日）。“已知系外行星数量几乎翻倍，美国宇航局发现715颗新世界”。space.com。检索日期：2014年2月27日。
\item 阿莫斯，乔纳森（2014年2月26日）。“开普勒望远镜发现大量行星”。BBC新闻。检索日期：2014年2月27日
\item 约翰逊，米歇尔；周，费利西亚（2015年7月23日）。“美国国家航空航天局开普勒任务发现地球的更大、更古老的表亲”. 美国国家航空航天局.
\item 美国国家航空航天局。“发现警报！奇特行星或将揭示其秘密”。系外行星探索：太阳系外的行星。检索日期：2018年11月28日
\item 富尔顿，本杰明·J.；佩蒂古拉，埃里克·A.；霍华德，安德鲁·W.；艾萨克森，霍华德；马西，杰弗里·W.；卡吉尔，菲利普·A.；赫布，莱斯利；魏斯，劳伦·M.；约翰逊，约翰·阿舍；莫顿，蒂莫西·D.；西努科夫，埃文；克罗斯菲尔德，伊恩·J. M.；希尔希，莉娅·A.（2017年9月1日）。“加利福尼亚-开普勒巡天研究Ⅲ：小型行星半径分布中的一个空白区*”. 天文期刊. 154 (3): 109. arXiv:1703.10375. 文献标识码:2017AJ....154..109F. DOI:10.3847/1538-3881/aa80eb. ISSN 0004-6256.
\item “半径间隙”.sites.astro.caltech.edu.检索日期2024年4月3日
\item 【视频】TOI 700d：一颗与地球大小相当的行星被发现在“宜居带”内.midilibre.fr（法语）。检索日期：2020年4月17日
\item “系外行星及候选体统计”。美国国家航空航天局系外行星档案，加州理工学院。检索日期：2020年1月17日。
\item 科伦，杰里（2013年11月4日）。“开普勒”。nasa.gov。美国国家航空航天局。于2013年11月5日从原文存档。检索日期：2013年11月4日
\item  哈灵顿，J. D.；约翰逊，M.（2013年11月4日）。“NASA开普勒成果开启天文学新时代”.
\item “NASA系外行星档案KOI表”。美国国家航空航天局。于2014年2月26日从原文存档。检索日期：2014年2月28日。
\item 刘文，莎拉（2017年6月19日）。“美国国家航空航天局开普勒太空望远镜发现数百颗新系外行星，使总数增至4,034颗”。美国国家航空航天局。检索日期：2017年6月19日
\item 奥布赖，丹尼斯（2017年6月19日）。“地球大小的行星跻身美国国家航空航天局开普勒望远镜最终统计结果”. 《纽约时报》。于2022年1月1日存档自原文。
\item 克雷恩，莉亚（2020年9月23日）。“天文学家可能发现了首个位于其他星系的行星”。《新科学家》。检索日期：2020年9月25日
\item 迪·斯塔法诺，R. 等人（2020年9月18日）。“M51-ULS-1b：首个外星系行星候选体”。arXiv:2009.08987 [astro-ph.HE].
\item 佩里曼，迈克尔（2011）。系外行星手册。剑桥大学出版社。第149. 。ISBN 978-0-521-76559-6.
\item 佩佩，F.；洛维斯，C.；塞格朗桑，D.；本茨，W.；布希，F.；杜穆斯克，X.；马约尔，M.；凯洛兹，D.；桑托斯，N. C.；于德里，S.（2011）。“HARPS在宜居带中寻找类地行星”。天文学与天体物理学. 534: A58。arXiv:1108.3447. 文献编码:2011A&A...534A..58P. DOI:10.1051/0004-6361/201117055. S2CID 15088852.
\item 行星探测：寻找类地行星 已归档2010-07-28 于互联网档案馆。科学计算，2010年7月19日
\item 巴拉德，S.；法布里基，D.；弗雷桑，F.；沙尔博诺，D.；德塞特，J. M.；托雷斯，G.；马西，G.；伯克，C. J.；艾萨克森，H.；亨泽，C.；施特芬，J. H.；恰尔迪，D. R.；霍威尔，S. B.；科克兰，W. D.；恩德尔，M.；布莱森，S. T.；罗伊，J. F.；霍尔曼，M. J.；利萨尔，J. J. 亨金斯，J. M. 亨金斯，M. 斯蒂尔，E. B. 福特，J. L. 克里斯蒂安森，C. K. 米多尔，M. R. 哈斯，J. 李，J. R. 霍尔，S. 麦考利夫，N. M. 巴塔利亚，D. G. 科赫等（2011）。“开普勒-19系统：一颗凌星的2.2倍地球半径行星及通过凌星 timing 变化探测到的第二颗行星”。天体物理学报.743(2): 200.arXiv:1109.1561Bibcode2011ApJ...743..200Bdoi10.1088/0004-637X/743/2/200S2CID42698813
\item 帕尔，A.；科奇斯，B.（2008）。“凌星系外行星系统中近日点进动的广义相对论水平测量”. 皇家天文学会月报. 389 (1): 191–198. arXiv:0806.0629. 文献标识码:2008MNRAS.389..191P. DOI:10.1111/j.1365-2966.2008.13512.x. S2CID 15282437.
\item 库里埃尔，萨尔瓦多；奥尔蒂斯-莱昂，吉塞拉·N.；米奥杜谢夫斯基，艾米·J.；桑切斯-贝尔穆德斯，乔尔（2022年9月）。“矮双星系统的三维轨道结构及其行星伴星”. 天文学杂志. 164 (3): 93. arXiv:2208.14553. 文献标识码:2022AJ....164...93C. DOI:10.3847/1538-3881/ac7c66. S2CID 251953478.
\item 索泽蒂，A.；皮纳蒙蒂，M. 等人（2023年9月）。“TNG上的GAPS计划。第47篇：谜题终获解答：HIP 66074b/Gaia-3b被确认为一颗质量巨大的巨行星，其轨道近乎正视且极度拉长”. 天文学与天体物理学. 677：L15。arXiv:2307.08653. 文献标识码:2023A&A...677L..15S. DOI:10.1051/0004-6361/202347329. HDL:2108/347124.
\item 达穆尔，蒂博（1992）。“相对论引力的强场检验与双脉冲星”. 物理评论D. 45 (6): 1840–1868. 文献标识码:1992PhRvD..45.1840D. DOI:10.1103/PhysRevD.45.1840. PMID 10014561.
\item 西尔沃蒂，R.；舒，S.；扬ulis，R.；索尔海姆，J.-E.；贝纳贝伊，S.；厄斯特森，R.；奥斯瓦尔特，T. D.；布鲁尼，I.；瓜兰迪，R.；博南诺，A.；沃克莱尔，G.；里德，M.；陈昌伟；莱博维茨，E.；帕帕罗，M.；巴兰，A.；沙尔皮内，S.；多莱兹，N.；卡瓦勒，S.；库尔茨，D.；莫斯卡利克，P.；里德尔，R.；佐拉，S.（2007）。“一颗围绕‘极端水平分支’恒星V 391 波江座运行的巨行星” （PDF）. 自然. 449 (7159): 189–191. 天体物理文献标识码:2007Natur.449..189S. DOI:10.1038/nature06143. PMID 17851517. S2CID 4342338.
\item 詹金斯，J. M.；多伊尔，劳伦斯·R.（2003年9月20日）。“利用空间光度计探测近距类木行星的反射光”。天体物理学杂志.1(595)：429–445。arXiv：astro-ph/0305473文献编码2003ApJ...595..429JDOI10.1086/377165S2CID17773111
\item 洛布，A.；高迪，B. S.（2003）。“由行星伴星引起的反射多普勒效应导致的恒星周期性流量变化”。天体物理学报快报. 588(2)：L117。arXiv:astro-ph/0303212. 文献标识码:2003ApJ...588L.117L. DOI:10.1086/375551. S2CID 10066891.
\item 阿特金森，南希（2013年5月13日）。“利用相对论和BEER寻找系外行星”。今日宇宙。检索日期：2023年2月12日
\item 施密德，H. M.；博济，J.-L.；费尔特，M.；吉斯勒，D.；格拉顿，R.；亨宁，T.；约斯，F.；卡斯珀，M.；伦岑，R.；穆耶，D.；穆图，C.；奎伦巴赫，A.；斯塔姆，D. M.；塔尔曼，C.；廷贝根，J.；韦里诺，C.；沃特斯，R.；沃尔斯滕克罗夫特，R.（2006）。“利用偏振测量法搜寻与研究太阳系外行星”. 国际天文学联合会会刊. 1: 165. 文献标识码:2006dies.conf..165S. DOI:10.1017/S1743921306009252.
\item 别尔久吉娜，S. V.；别尔久金，A. V.；弗卢里，D. M.；皮罗拉，V.（2008）。“首次探测到来自系外行星大气的偏振散射光”。天体物理学报. 673(1)：L83。arXiv:0712.0193. Bibcode:2008ApJ...673L..83B. doi:10.1086/527320. S2CID 14366978.
\item 丹杰洛，G.；杜里森，R. H.；利绍尔，J. J. （2011）。“巨行星的形成”。载于 S. 西格尔（编）。系外行星。亚利桑那大学出版社，图森，亚利桑那州。第319–346. arXiv:1006.5486. Bibcode:2010exop.book..319D.
\item 丹杰洛，G.；利绍尔，J. J.（2018）。“巨行星的形成”。载于迪格 H.、贝尔蒙特 J. 主编，《系外行星手册》。施普林格国际出版股份公司，施普林格自然集团成员。第2319–2343页。arXiv:1806.05649。文献编码2018haex.bookE.140DDOI10.1007/978-3-319-55333-7_140ISBN978-3-319-55332-0S2CID116913980
\item 卡尔韦特，努丽亚；达莱西奥，保拉；哈特曼，李；威尔纳，大卫；沃尔什，安德鲁；西特科，迈克尔（2001）。“1000万年老原行星盘中正在形成的空隙的证据”。天体物理学报. 568 (2): 1008–1016. arXiv:astro-ph/0201425. 文献标识码:2002ApJ...568.1008C. DOI:10.1086/339061. S2CID 8706944.
\item 弗里德伦德，马尔科姆；盖多斯，埃里克；巴拉甘，奥斯卡；佩尔松，卡丽娜；甘多尔菲，达维德；卡布雷拉，胡安；平野照之；久保原正幸；奇兹马迪亚，Sz；诺瓦克，格热戈日；恩德尔，迈克尔；格尔齐瓦，萨沙；科尔特，朱迪思；普法夫，耶雷米亚斯；比奇，贝特拉姆；约翰森，安德斯；穆斯蒂尔，亚历山大；戴维斯，梅尔文；迪格，汉斯；帕耶，恩里克；科克兰，威廉；艾格米勒，菲利普；埃里克森，安德斯；根特，艾克；哈策斯，阿蒂；基勒里希，阿曼达；久藤，智之；麦克奎因，菲利普；成田，典雄；内斯普拉尔，大卫；佩茨奥尔德，马丁；普列托-阿兰斯，豪尔赫；劳尔，海克；范艾伦，文森特（2017年4月28日）。“EPIC210894022b——一颗短周期超地球，环绕着一颗金属含量低、已演化的古老恒星运行”。天文学与天体物理学. 604: A16。arXiv:1704.08284. doi:10.1051/0004-6361/201730822. S2CID 39412906.
\item 丹杰洛，G.；博登海默，P.（2016）。“开普勒11行星的原位与异地形成模型”. 天体物理学杂志. 828（1）：第33页（共32页）。arXiv:1606.08088. 文献标识码:2016ApJ...828...33D. DOI:10.3847/0004-637X/828/1/33. S2CID 119203398.
\item 丹杰洛，G.；博登海默，P.（2013）。“嵌入原行星盘中的年轻行星包层的三维辐射流体动力学计算”。天体物理学杂志. 778(1)：77（29页）。arXiv:1310.2211. 文献标识码:2013ApJ...778...77D. DOI:10.1088/0004-637X/778/1/77. S2CID 118522228.
\item 丹杰洛，G.；魏登席林，S. J.；利绍尔，J. J.；博登海默，P.（2014）。“木星的形成：富含低质量包层对核心吸积的增强作用”。国际行星科学杂志. 241: 298–312. arXiv:1405.7305. 文献标识码:2014Icar..241..298D. DOI:10.1016/j.icarus.2014.06.029. S2CID 118572605.
\item Lammer, H.; Stokl, A.; Erkaev, N. V.; Dorfi, E. A.; Odert, P.; Gudel, M.; Kulikov, Y. N.; Kislyakova, K. G.; Leitzinger, M. (2014). "Origin and loss of nebula-captured hydrogen envelopes from 'sub'- to 'super-Earths' in the habitable zone of Sun-like stars". Monthly Notices of the Royal Astronomical Society. 439 (4): 3225–3238. arXiv:1401.2765. Bibcode:2014MNRAS.439.3225L. doi:10.1093/mnras/stu085. S2CID 118620603.
\item 约翰逊，R. E.（2010）。“热驱动的大气逃逸”。天体物理学杂志. 716 (2): 1573–1578. arXiv:1001.0917. 文献编码:2010ApJ...716.1573J. DOI:10.1088/0004-637X/716/2/1573. S2CID 36285464.
\item 赞德哈斯，J.；塞古拉，A.；拉加，A.C.（2010）。“主序M型恒星周围行星的恒星风导致的大气质量损失”。国际行星科学杂志.210(2)：539–544。arXiv：1006.0021天文学文献标识码2010Icar..210..539ZDOI10.1016/j.icarus.2010.07.013S2CID119243879
\item 增田，K.（2014）。“利用凌星 timing 变化揭示开普勒-51周围极低密度行星，以及一种类似行星间食现象的异常现象”。天体物理学报. 783 (1): 53. arXiv:1401.2885. Bibcode:2014ApJ...783...53M. doi:10.1088/0004-637X/783/1/53. S2CID 119106865.
\item “艺术家对围绕两颗恒星运行的系外行星的想象图”.www.spacetelescope.org。检索日期：2016年9月24日
\item 卡桑，A.；库巴斯，D.；博利厄，J.-P.；多米尼克，M.；霍恩，K.；格林希尔，J.；瓦姆布桑斯，J.；门齐斯，J.；威廉姆斯，A.；约尔根森，U. G.；乌达尔斯基，A.；贝内特，D. P.；阿尔布罗，M. D.；巴蒂斯塔，V.；布里扬，S.；卡尔德威尔，J. A. R.；科尔，A.；库图尔，C.；库克，K. H.；迪特斯，S.；普雷斯特，D. D.；多纳托维奇，J.；福克，P.；希尔，K.；凯恩斯，N.；凯恩，S.；马尔凯特，J.-B.；马丁，R.；波拉德，K. R.；萨胡，K. C.（2012年1月11日）。“通过微引力透镜观测发现，银河系每颗恒星至少有一颗或更多被束缚的行星。”自然. 481 (7380): 167–169. arXiv:1202.0903. Bibcode:2012Natur.481..167C. doi:10.1038/nature10684. PMID 22237108. S2CID 2614136.
\item 佩蒂古拉，E. A.；霍华德，A. W.；马西，G. W.（2013）。“围绕类太阳恒星运行的地球大小行星的普遍性”. 美国国家科学院院刊. 110 (48): 19273–19278. arXiv:1311.6806. 天文学文献标识码:2013PNAS..11019273P. DOI:10.1073/pnas.1319909110. PMC 3845182. PMID 24191033.
\item 卡明，安德鲁；巴特勒，R·保罗；马西，杰弗里·W.; 沃格特，史蒂文·S.；赖特，杰森·T.；费舍尔，黛布拉·A.（2008）。“凯克行星搜索：可探测性及系外行星的最小质量和轨道周期分布”。太平洋天文学会出版物. 120 (867): 531–554. arXiv:0803.3357. Bibcode:2008PASP..120..531C. doi:10.1086/588487. S2CID 10979195.
\item 邦菲尔，泽维尔；福尔韦耶，蒂埃里；德尔福斯，泽维尔；于德里，斯蒂芬；马约尔，米歇尔；佩里耶，克里斯蒂安；布希，弗朗索瓦；佩佩，弗朗切斯科；凯洛，迪迪埃；贝尔托，让-卢普（2005）。“HARPS对南天系外行星的搜寻VI：一颗质量与海王星相当的行星围绕着邻近的M型矮星格利泽581。”天文学与天体物理学. 443 (3): L15 – L18. arXiv:astro-ph/0509211. 文献标识码:2005A&A...443L..15B. DOI:10.1051/0004-6361:200500193. S2CID 59569803.
\item "Kepler / K2 - NASA Science". 23 May 2023. Retrieved 22 August 2025.
\item 王，J.；费舍尔，D. A.（2014）。“揭示不同太阳型恒星行星的普遍行星-金属丰度相关性”。天文学期刊. 149 (1): 14. arXiv:1310.7830. Bibcode:2015AJ....149...14W. doi:10.1088/0004-6256/149/1/14. S2CID 118415186.
\item “科学工作”.www.univie.ac.at. 检索日期2022年1月17日
\item “STAR-DATA”.www.univie.ac.at。检索日期：2022年1月17日
\item 系外行星及其恒星的视亮度与大小存档于2014年8月12日，互联网档案馆，阿贝尔·门德斯，更新于2012年6月30日，下午12:10
\item 加纳，罗布（2016年10月31日）。“NASA哈勃望远镜发现一颗真正的蓝色行星”。NASA。检索日期：2022年1月17日
\item 埃文斯，T. M.；蓬特，F. D. R.；辛格，D. K.；艾格兰，S.；巴尔斯托，J. K.；德塞，J. M.；吉布森，N.；恒，K.；克努特森，H. A.；勒卡韦耶·德斯昂日，A.（2013）。“HD189733b的深蓝色：哈勃空间望远镜/空间望远镜成像光谱仪在可见波长下的反照率测量”。天体物理学报. 772(2): L16.arXiv:1307.3239. Bibcode:2013ApJ...772L..16E. doi:10.1088/2041-8205/772/2/L16. S2CID 38344760.
\item 久原，M.；田村，M.；工藤，T.；扬森，M.；神取，R.；布兰特，T. D.；塔尔曼，C.；施皮格尔，D.；比勒，B.；卡森，J.；堀，Y.；铃木，R.；伯罗斯，亚当；亨宁，T.；特纳，E. L.；麦克埃尔温，M. W.；莫罗-马丁，A.；末永，T.；高桥，Y. H.权，J.卢卡斯，L.阿贝，W.布兰德纳，S.埃格纳，M.费尔特，H.藤原，M.后藤，C.A.格雷迪，O.居永，J.桥本等（2013）。“围绕类太阳恒星GJ 504运行的寒冷木星型系外行星的直接成像” （PDF）. 天体物理学报. 774 (11): 11. arXiv:1307.2886. Bibcode:2013ApJ...774...11K. doi:10.1088/0004-637X/774/1/11. S2CID 53343537.
\item 卡森；塔尔曼；扬森；科扎基斯；博纳福伊；比勒；施利德；柯里；麦克埃尔温（2012年11月15日）。“在晚期B型恒星卡帕人马座周围直接成像发现一颗‘超级木星’”。天体物理学杂志. 763（2）：L32。arXiv:1211.3744. 文献编码:2013ApJ...763L..32C. DOI:10.1088/2041-8205/763/2/L32. S2CID 119253577.
\item “被氦气包裹的行星在我们的银河系中可能很常见”。SpaceDaily，2015年6月16日。检索日期：2015年8月3日。
\item “煤炭般漆黑的外星行星是有史以来最暗的天体”。Space.com，2011年8月11日。检索日期：2011年8月12日。
\item 基平，大卫·M.；施皮格尔，大卫·S.（2011）。“从最黑暗的世界探测可见光”. 皇家天文学会月报：快报. 417 (1): L88－ L92. arXiv:1108.2297. Bibcode:2011MNRAS.417L..88K. doi:10.1111/j.1745-3933.2011.01127.x. S2CID 119287494.
\item 巴克莱，T.；胡贝尔，D.；罗伊，J. F.；福特尼，J. J.；莫利，C. V.；昆塔纳，E. V.；法布里基，D. C.；巴伦森，G.；布洛门，S.；克里斯蒂安森，J. L.；德莫里，B. O.；富尔顿，B. J.；詹金斯，J. M.；穆拉利，F.；拉戈齐内，D.；希德，S. E.；什波雷尔，A.；特嫩鲍姆，P.；汤普森，S. E.（2012）。“通过测光法测定TrES-2系统中行星与恒星的质量和半径”。天体物理学杂志. 761 (1): 53. arXiv:1210.4592. 文献标识码:2012ApJ...761...53B. DOI:10.1088/0004-637X/761/1/53. S2CID 18216065.
\item 伯罗兹，亚当（2014）。《利用WFIRST进行日冕仪系外行星成像与光谱学的科学回报》。arXiv:1412.6097 [astro-ph.EP].
\item 查尔斯·Q·崔（2014年11月20日）。“揭开外星世界磁场的秘密”。太空网。检索日期：2022年1月17日
\item 基斯利亚科娃，K. G.；霍尔姆斯特伦，M.；拉默，H.；奥德特，P.；霍达琴科，M. L.（2014）。“基于莱曼α观测确定的HD 209458b的磁矩与等离子体环境”。科学. 346 (6212): 981–984. arXiv:1411.6875. Bibcode:2014Sci...346..981K. doi:10.1126/science.1257829. PMID 25414310. S2CID 206560188.
\item 尼科尔斯，J. D.（2011）。“具有内部等离子体源的木星型系外行星中的磁层-电离层耦合：对极光射电辐射可探测性的启示”. 皇家天文学会月报. 414 (3): 2125–2138. arXiv:1102.2737. 文献标识码:2011MNRAS.414.2125N. DOI:10.1111/j.1365-2966.2011.18528.x. S2CID 56567587.
\item “射电望远镜或有助于发现系外行星”.Redorbit，2011年4月18日，检索日期：2022年1月17日
\item “系外行星的射电探测：现状与未来展望”（PDF）。美国海军研究实验室、NASA/戈达德太空飞行中心、美国国家射电天文台、巴黎天文台。于2008年10月30日从原始网页存档。检索日期：2008年10月15日
\item 路线，马修（2024年5月1日）。“罗马。IV. 在5 GHz频段对据称存在系外行星的系统中进行亚恒星磁层射电辐射的阿雷西博搜寻”。天体物理学报966（1）：55。arXiv：2403.02226文献标识码2024ApJ...966...55Rdoi10.3847/1538-4357/ad30ff
\item 基恩，萨姆（2016）。“禁用植物，禁用化学”。蒸馏2（2）：5。于2018年3月23日从原文存档。检索日期：2018年3月22日
\item 崔，查尔斯·Q.（2012年11月22日）。“超级地球获得由液态金属提供的磁性‘防护罩’”。Space.com。检索日期：2022年1月17日
\item 布扎西，D.（2013）。“恒星磁场作为系外巨行星的加热源”。天体物理学杂志. 765（2）：L25。arXiv:1302.1466. 文献标识码:2013ApJ...765L..25B. doi:10.1088/2041-8205/765/2/L25. S2CID 118978422.
\item 张健，2018年8月16日。“利用168台巨型激光器化解关于氢的争议——劳伦斯利弗莫尔国家实验室的科学家表示，他们在一项旨在理解液态金属氢状态的实验中‘正逐步接近真相’”。纽约时报。于2022年1月1日存档自原文。检索日期：2018年8月18日
\item 工作人员（2018年8月16日）。“在高压下，氢展现出对巨行星内部结构的反映——氢是宇宙中含量最丰富的元素，也是最简单的元素，但这种简单却具有欺骗性。”《科学日报》。检索日期：2018年8月18日
\item 罗特，马修（2019年2月10日）。“罗马的崛起。I. 对HD 189733系统中恒星与行星相互作用的多波长分析”。《天体物理学杂志》872（1）：79。arXiv：1901.02048文献标识码2019ApJ...872...79RDOI10.3847/1538-4357/aafc25S2CID119350145
\item 路线，马修；沃尔什赞，亚历克斯（2023年8月1日）。“罗马：III. 阿雷西博5 GHz频段恒星-行星相互作用搜寻”. 天体物理学报. 952 (2): 118. arXiv:2202.08899. 文献标识码:2023ApJ...952..118R. doi:10.3847/1538-4357/acd9ad.
\item 拉比，帕桑特（2019年7月29日）。“‘热木星’系外行星的磁场比我们想象的要强得多”。太空网。检索日期：2022年1月17日
\item 考利，P·威尔逊；施科尔尼克，叶夫根娅·L.；拉马，乔；兰扎，安东尼诺·F.（2019年12月）。“来自恒星与行星相互作用信号的热木星磁场强度”。自然天文. 3 (12): 1128–1134. arXiv:1907.09068. Bibcode:2019NatAs...3.1128C. doi:10.1038/s41550-019-0840-x. ISSN 2397-3366. S2CID 198147426.
\item 奥卡拉汉，乔纳森（2023年8月7日）。“系外行星或可帮助我们了解行星如何产生磁性”. 量子杂志.
\item 瓦伦西亚，黛安娜；奥康奈尔，理查德·J.（2009）。“地球与超级地球上的对流尺度与俯冲作用”。地球与行星科学快报. 286 (3–4): 492–502. 文献标识码:2009E&PSL.286..492V. DOI:10.1016/j.epsl.2009.07.015.
\item 范赫克，H.J.；塔克利，P.J.（2011）。“超级地球上的板块构造：与地球一样甚至更有可能。”地球与行星科学快报. 310 (3–4): 252–261. 文献标识码:2011E&PSL.310..252V. DOI:10.1016/j.epsl.2011.07.029.
\item 奥尼尔，C.；莱纳迪奇，A.（2007）。“超大地球的地质后果”. 地球物理研究快报. 34（19）：L19204。文献标识码:2007GeoRL..3419204O. DOI:10.1029/2007GL030598. S2CID 41617531.
\item 瓦伦西亚，黛安娜；奥康奈尔，理查德·J.；萨塞洛夫，迪米塔尔·D（2007年11月）。“超级地球上板块构造的必然性”。天体物理学报快报.670(1)：L45－L48arXiv：0710.0699Bibcode2007ApJ...670L..45Vdoi10.1086/524012S2CID9432267
\item “超级地球很可能同时拥有海洋和大陆——天体生物学”.astrobiology.com. 2014年1月7日. 检索日期：2022年1月17日
\item 考万，N. B.；阿博特，D. S.（2014）。“海洋与地幔之间的水循环：超级地球未必是‘水世界’”。天体物理学杂志. 781 (1): 27. arXiv:1401.0720. Bibcode:2014ApJ...781...27C. doi:10.1088/0004-637X/781/1/27. S2CID 56272100.
\item 莱莫尼克，迈克尔·D.（2015年5月6日）。“天文学家可能已在距离地球40光年的地方发现了火山”。国家地理。于2015年5月9日从原文存档。检索日期：2015年11月8日
\item 德莫里，布里斯-奥利维耶；吉隆，迈克尔；马杜苏丹，尼库；凯洛兹，迪迪埃（2015）。“超地球55 Cnc e的变异性”. 皇家天文学会月报. 455 (2): 2018–2027. arXiv:1505.00269. 文献标识码:2016MNRAS.455.2018D. DOI:10.1093/mnras/stv2239. S2CID 53662519.
\item {{引用网络 |url=https://www.sci.news/astronomy/volcano-covered-exo-earth-lp-791-18d-11930.html|标题=发现一颗被火山覆盖的类地系外行星，环绕着附近恒星运行 |网站=Sci.News}}
\item 贝洛-阿鲁费，亚伦；达米亚诺，马里奥；等（2025年1月）。“类地行星L98-59b存在火山大气的证据”。天体物理学报快报980（2）：L26。arXiv：2501.18680文献标识码2025ApJ...980L..26Bdoi10.3847/2041-8213/adaf22
\item “科学家发现一个类土星环系统正遮蔽一颗类太阳恒星”.www.spacedaily.com。检索日期：2022年1月17日
\item 马马杰克，E. E.；奎伦，A. C.；佩考，M. J.；穆勒坎普，F.；斯科特，E. L.；肯沃西，M. A.；卡梅隆，A. C.；帕利，N. R.（2012）。“掩食中的行星形成区：发现一个绕行年轻类太阳恒星运行的系外环系统，以及未来探测次级盘和行星盘遮蔽现象的前景。”天文学杂志. 143 (3): 72. arXiv:1108.4070. Bibcode:2012AJ....143...72M. doi:10.1088/0004-6256/143/3/72. S2CID 55818711.
\item 里德，史蒂文；肯沃西，马修·A.（2016年11月）。“J1407b环系大小与动力学的限制”. 天文学与天体物理学. 596: 5. arXiv:1609.08485. 文献标识码:2016A&A...596A...9R. doi:10.1051/0004-6361/201629567. S2CID 118413749。A9.
\item 肯沃西，M. A.；克拉森，P. D.；闵，M.；范德马雷尔，N.；博恩，A. J.；卡玛，M. 等人（2020年1月）。“针对年轻系外行星凌星系统J1407（半人马座V1400）的ALMA与NACO观测”. 天文学与天体物理学. 633: 6. arXiv:1912.03314. Bibcode:2020A&A...633A.115K. doi:10.1051/0004-6361/201936141. A115.
\item 萨伊恩费斯特，M.；苏利斯，S.；沙尔庞蒂耶，P.；桑特纳，A.（2023）。“迁移中的系外卫星形成的斜环：一种可能的长周期系外行星半径增大的起源”。天文学与天体物理学.675DOI：10.1051/0004-6361/20234674（已于2025年7月12日失效）。：CS1维护：截至2025年7月，DOI无效（链接)
\item 阿拉姆，穆纳扎·K.；柯克，詹姆斯；德雷辛，考特妮·D.；洛佩斯-莫拉莱斯，梅赛德斯；大野和正；高鹏；阿金萨尼，巴巴图德；桑特尔内，亚历山大；格鲁法尔，萨洛梅；阿迪贝基扬，瓦尔丹；巴罗斯，苏珊娜·C. C.；布赫哈维，拉斯·A.；克罗斯菲尔德，伊恩·J. M.；戴，飞；德勒伊，玛格丽；贾卡隆，史蒂文；利略-博克斯，豪尔赫；马利，马克；梅奥，安德鲁·W.；莫蒂埃，安内莉丝；桑托斯，努诺·C.；索萨，塞尔吉奥·G.；图特尔布姆，艾玛·V.；惠特利，彼得·J.；范德伯格，安德鲁·M.（2022）。“HIP 41378 f 的首张近红外透射光谱：一个多行星系统中的一颗低质量、适居的木星型行星”. 天体物理学报快报. 927(1): L5.arXiv:2201.02686. Bibcode:2022ApJ...927L...5A. doi:10.3847/2041-8213/ac559d. S2CID 245837282.
\item 施利希廷，希尔克·E.；张菲利普（2011）。“温暖的土星：关于位于冰线以内系外行星周围环状结构的本质”。天体物理学报. 734 (2): 117. arXiv:1104.3863. 文献编码:2011ApJ...734..117S. doi:10.1088/0004-637X/734/2/117. S2CID 42698264.
\item 俄罗斯天文学家首次在系外行星附近发现卫星（俄文）——“对WASP-12b亮度变化曲线的研究为俄罗斯天文学家带来了非同寻常的发现：他们发现了规律性的闪光现象……尽管恒星表面的黑子也会引起类似的亮度变化，但观测到的闪光在持续时间、波形和振幅上都非常相似，这有力地证明了系外卫星存在的可能性。”
\item 基斯利亚科娃，K. G.；皮拉特-洛辛格，E.；冯克，B.；拉默，H.；福萨蒂，L.；埃格尔，S.；施瓦茨，R.；布杰达，M. Y.；埃尔卡耶夫，N. V.（2016），“关于WASP-12和HD 189733系统的紫外异常：特洛伊卫星作为等离子体源”，皇家天文学会月报, 461 (1): 988–999, arXiv:1605.02507, Bibcode:2016MNRAS.461..988K, doi:10.1093/mnras/stw1110, S2CID 119205132
\item 贝内特，D. P.；巴蒂斯塔，V.；邦德，I. A.；贝内特，C. S.；铃木，D.；博多利耶，J. -P.；乌达尔斯基，A.；多纳托维奇，J.；博扎，V.；阿部，F.；博茨勒，C. S.；弗里曼，M.；福冈，D.；福井，A.；伊藤，Y.；小野，N.；林，C. H.；增田，K.；松原，Y.；村木洋；南波诚；大西健；拉滕伯里N.J.；齐藤隆；沙利文D.J.；住米隆；斯威特曼W.L.；特里斯特拉姆P.J.；津留美直树；和田康；等（2014）。“MOA-2011-BLG-262Lb：一颗亚地球质量的卫星，环绕着一颗气态巨行星运行，或位于银河系核球中的高速行星系统。”天体物理学报. 785 (2): 155. arXiv:1312.3951. Bibcode:2014ApJ...785..155B. doi:10.1088/0004-637X/785/2/155. S2CID 118327512.
\item 特里，肖恩·K.；博利厄，让-菲利普；贝内特，大卫·P.；巴塔查里亚，阿帕尔纳；胡尔伯格，乔恩；赫斯顿，梅西·J.；小越本直树；布莱克曼，乔舒亚·W.；邦德，伊恩·A.（2025）。“银河系核球中一颗候选高速系外行星系统”. 天文学杂志. 169 (3): 131. arXiv:2410.09147. Bibcode:2025AJ....169..131T. doi:10.3847/1538-3881/ad9b0f.
\item 蒂奇，亚历克斯；基平，大卫·M.（2018年10月1日）。“一颗围绕开普勒-1625b运行的大型系外卫星的证据”. 科学进展. 4（10）eaav1784。arXiv:1810.02362. 文献编码:2018SciA....4.1784T. DOI:10.1126/sciadv.aav1784. ISSN 2375-2548. PMC 6170104. PMID 30306135.
\item “科学家认为他们在遥远的恒星系统中发现了一颗巨大而奇特的卫星”.NPR.org。检索日期：2022年3月28日
\item 赫勒，雷内；希普克，迈克尔（2024）。“开普勒-1625 b和开普勒-1708 B周围不太可能存在大型系外卫星”. 自然天文. 8 (2): 193. arXiv:2312.03786. 文献标识码:2024NatAs...8..193H. DOI:10.1038/s41550-023-02148-w.
\item 基平，大卫；蒂奇，亚历克斯；亚哈洛米，丹尼尔·A.；卡塞塞，本；夸尔斯，比利；布莱森，史蒂夫；汉森，布拉德；苏拉吉，朱迪特；伯克，克里（2024年1月18日）。“对：开普勒-1625 b和开普勒-1708 b周围不太可能存在大型系外卫星”的回应”。arXiv:2401.10333 [astro-ph.EP].
\item 奥扎，阿普尔瓦·V.；约翰逊，罗伯特·E.；勒卢什，埃曼努埃尔；施密特，卡尔；施奈德，尼克；黄晨亮；甘博里诺，黛安娜；格贝克，安德烈娅；维滕巴赫，奥雷利安；德莫里，布里斯-奥利维耶；莫尔达西尼，克里斯托夫；萨克赛纳，普拉巴尔；杜布瓦，大卫；穆莱，阿里埃尔；托马斯，尼古拉斯（2019年8月28日）。“围绕近距离气态巨行星运行的火山卫星中的钠和钾特征”。天体物理学杂志885（2）：168。arXiv：1908.10732文献编码2019ApJ...885..168ODOI10.3847/1538-4357/ab40ccS2CID201651224
\item 格贝克，安德烈娅；奥扎，阿普尔瓦（2020年7月29日）。“系外行星系统的碱性外层大气：蒸发传输光谱”。《皇家天文学会月报》497（4）：5271–5291。arXiv：2005.02536文献标识码2020MNRAS.497.5271GDOI10.1093/mnras/staa2193S2CID218516741。检索于2020年12月8日
\item 两颗系外行星上的阴天与晴天大气.www.spacetelescope.org. 检索日期：2017年6月6日
\item 夏尔博诺，大卫；等（2002）。“系外行星大气层的探测”。天体物理学报. 568 (1): 377–384. arXiv:astro-ph/0111544. 文献标识码:2002ApJ...568..377C. DOI:10.1086/338770. S2CID 14487268.
\item 马杜苏丹，尼库；克努特森，希瑟；福特尼，乔纳森；巴尔曼，特拉维斯（2014）。“系外行星大气层”。原恒星与行星VI。第739页。arXiv:1402.1169. Bibcode:2014prpl.conf..739M. doi:10.2458/azu_uapress_9780816531240-ch032. ISBN 978-0-8165-3124-0. S2CID 118337613.
\item 西格尔，S.；德明，D.（2010）。“系外行星大气层”。天文学与天体物理学年评. 48: 631–672. arXiv:1005.4037. Bibcode:2010ARA&A..48..631S. doi:10.1146/annurev-astro-081309-130837. S2CID 119269678.
\item 罗德勒，F.；洛佩斯-莫拉莱斯，M.；里巴斯，I.（2012年7月）。“对非凌星热木星τ宝瓶座b的质量测量”。天体物理学报快报. 753(1)：L25。arXiv:1206.6197. 文献编码:2012ApJ...753L..25R. DOI:10.1088/2041-8205/753/1/L25. S2CID 119177983。L25.
\item 布罗吉，M.；斯内伦，I. A. G.；德科克，R. J.；阿尔布雷希特，S.；伯克比，J.；德穆伊，E. J. W.（2012）。“来自τ牧夫座b行星昼侧的轨道运动特征”。自然. 486 (7404): 502–504. arXiv:1206.6109. Bibcode:2012Natur.486..502B. doi:10.1038/nature11161. PMID 22739313. S2CID 4368217.
\item 圣弗勒，尼古拉斯（2017年5月19日）。“从百万英里之外观测地球上神秘的闪光”。《纽约时报》。于2022年1月1日存档自原文。检索日期：2017年5月20日
\item 马尔沙克，亚历山大；瓦尔奈，塔马什；科斯京斯基，亚历山大（2017年5月15日）。“从深空观测到的陆地反照：在拉格朗日点探测到定向冰晶”. 地球物理研究快报. 44 (10): 5197–5202. 文献编码:2017GeoRL..44.5197M. DOI:10.1002/2017GL073248. HDL:11603/13118. S2CID 109930589.
\item 莱顿大学。“蒸发的系外行星掀起尘埃”。phys.org。检索日期：2022年1月17日
\item “新发现的系外行星正在蒸发殆尽”。TGDaily，2012年5月18日。检索日期：2022年1月17日
\item 巴努，辛迪娅·N.（2015年6月25日）。“一颗拥有九百万英里长尾巴的行星”。《纽约时报》。检索日期：2015年6月25日
\item 雷蒙德，肖恩（2015年2月20日）。“别再提‘类地行星’了——我们首先会在眼球行星上发现外星生命”。Nautilus。于2017年6月23日存档自原文。检索日期：2022年1月17日
\item 多布罗沃尔斯基，安东尼·R.（2015）。“偏心系外行星上的日照模式”。土星. 250: 395–399. 文献标识码:2015Icar..250..395D. DOI:10.1016/j.icarus.2014.12.017.
\item 多布罗沃尔斯基斯，安东尼·R.（2013）。“具有偏心率和倾角的系外行星上的日照”。国际行星科学杂志. 226 (1): 760–776. 文献标识码:2013Icar..226..760D. doi:10.1016/j.icarus.2013.06.026.
\item 胡仁宇；埃赫尔曼，贝瑟尼·L.；西格尔，萨拉（2012）。“类地系外行星表面的理论光谱”。天体物理学报. 752 (1): 7. arXiv:1204.1544. Bibcode:2012ApJ...752....7H. doi:10.1088/0004-637X/752/1/7. S2CID 15219541.
\item “美国国家航空航天局、欧洲航天局以及K. 海恩斯和A. 曼德尔（戈达德太空飞行中心）”。检索日期：2015年6月15日。
\item 克努特森，H. A.；沙尔博诺，D.；艾伦，L. E.；福特尼，J. J.；阿戈尔，E.；考万，N. B.；肖曼，A. P.；库珀，C. S.；梅吉斯，S. T.（2007）。“系外行星HD 189733b昼夜对比图” （PDF）. 自然. 447 (7141): 183–186. arXiv:0705.0993. 文献标识码:2007Natur.447..183K. DOI:10.1038/nature05782. PMID 17495920. S2CID 4402268.
\item 奥利维耶，马克；莫雷尔，玛丽-克里斯汀（2014）。“行星环境与生命起源”. 生物会议网络. 2: 00001. DOI:10.1051/bioconf/20140200001.
\item “氧气并非系外行星上存在生命的决定性证据”。国立天文台。天体生物学网站。2015年9月10日。检索日期2015年9月11日
\item 科普拉普，拉维·库马尔（2013）。“对开普勒M型矮星宜居带内类地行星出现率的修订估算”。天体物理学报快报. 767(1)：L8。arXiv:1303.2649. 文献编码:2013ApJ...767L...8K. doi:10.1088/2041-8205/767/1/L8. S2CID 119103101.
\item 克鲁兹，玛丽亚；库恩茨，罗伯特（2013）。“系外行星——专题介绍”. 科学. 340 (6132): 565. DOI:10.1126/science.340.6132.565. PMID 23641107.
\item 崔，查尔斯·Q.（2011年9月1日）。“‘沙丘’行星上存在外星生命的可能性更大”. 天体生物学杂志。于2013年12月2日从原文存档。
\item 阿部洋；阿部大介；斯利普，N. H.；扎恩勒，K. J.（2011）。“干燥行星的宜居带边界”。天体生物学. 11 (5): 443–460. 文献标识码:2011AsBio..11..443A. DOI:10.1089/ast.2010.0545. PMID 21707386.
\item 西格尔，S.（2013）。“系外行星宜居性”。科学. 340 (6132): 577–581. 文献标识码:2013Sci...340..577S. CiteSeerX 10.1.1.402.2983. DOI:10.1126/science.1232226. PMID 23641111. S2CID 206546351.
\item 科普拉普，拉维·库马尔；拉米雷斯，拉姆塞斯·M.；肖特尔科特，詹姆斯；卡斯廷，詹姆斯·F.；多马加尔-戈德曼，肖恩；艾梅，文森特（2014）。“主序星周围宜居带：与行星质量的依赖关系”。天体物理学杂志. 787（2）：L29。arXiv:1404.5292. 文献标识码:2014ApJ...787L..29K. DOI:10.1088/2041-8205/787/2/L29. S2CID 118588898.
\item 滨野克；阿部洋；玄田宏（2013）。“岩浆海洋固化时两类类地行星的形成”。自然. 497 (7451): 607–610. 文献标识码:2013Natur.497..607H. DOI:10.1038/nature12163. PMID 23719462. S2CID 4416458.
\item 杨，J.；布埃，G. L.；法布里基，D. C.；阿博特，D. S.（2014）。“宜居带内边缘对行星自转速率的强烈依赖性”（PDF）。天体物理学报787（1）：L2。arXiv：1404.4992文献编码2014ApJ...787L...2YDOI10.1088/2041-8205/787/1/L2S2CID56145598。于2016年4月12日从原始网页存档。检索日期：2016年7月28日
\item “现实版科幻世界 #2：炽热眼球行星”. planetplanet，2014年10月7日。
\item 杨俊；科万，尼古拉斯·B.；阿博特，多里安·S.（2013）。“稳定云反馈可显著扩大潮汐锁定行星的宜居带。”天体物理学报. 771(2)：L45。arXiv:1307.0515. 文献编码:2013ApJ...771L..45Y. DOI:10.1088/2041-8205/771/2/L45. S2CID 14119086.
\item 斯塔尔，米歇尔（2023年4月19日）。“科学家认为他们已缩小了最有可能孕育生命的恒星系统范围”。ScienceAlert。检索日期：2023年4月19日
\item 沙皮罗，安娜·V.；等（2023年4月18日）。“富含金属的恒星不太适合其行星上生命的演化”. 自然通讯. 14 (1893) 1893. 文献标识码:2023NatCo..14.1893S. DOI:10.1038/s41467-023-37195-4. PMC 10113254. PMID 37072387.
\item 阿蒙德，J. P.；泰斯克，A.（2005）。“深部地下微生物学前沿的拓展”。古地理学、古气候学、古生态学. 219 (1–2): 131–155. 文献标识码:2005PPP...219..131A. doi:10.1016/j.palaeo.2004.10.018.
\item 研究人员称，更遥远的行星“可能支持生命”。BBC新闻，2014年1月7日，检索日期：2023年2月12日
\item 阿博特，D. S.；斯威策，E. R.（2011）。“草原之狼：一颗宜居行星在恒星际空间中的提案”. 天体物理学杂志. 735（2）：L27。arXiv:1102.1108. 文献编码:2011ApJ...735L..27A. DOI:10.1088/2041-8205/735/2/L27. S2CID 73631942.
\item 洛布，A.（2014）。“早期宇宙的宜居时期”。国际天体生物学杂志。13(4)：337–339。arXiv：1312.0613Bibcode2014IJAsB..13..337LCiteSeerX10.1.1.748.4820doi10.1017/S1473550414000196S2CID2777386
\item 里奇韦，安迪（2015年7月29日）。“家园，甜蜜的系外卫星：寻找外星生命的全新前沿”。《新科学家》。检索日期：2023年2月12日
\item 林森迈尔，曼努埃尔；帕斯卡莱，萨尔瓦托雷；卢卡里尼，瓦莱里奥（2014）。“具有高倾角和偏心轨道的类地行星的宜居性：基于一般环流模型的研究结果”。行星与空间科学. 105: 43–59. arXiv:1401.5323. 文献标识码:2015P&SS..105...43L. DOI:10.1016/j.pss.2014.11.003. S2CID 119202437.
\item 凯利，彼得（2014年4月15日）。“天文学家：‘倾斜世界’可能孕育生命”。华盛顿大学新闻。检索日期：2023年2月12日
\item 阿姆斯特朗，J. C.；巴恩斯，R.；多马加尔-戈德曼，S.；布赖纳，J.；奎因，T. R.；米德斯，V. S.（2014）。“极端倾角变化对系外行星宜居性的影响”. 天体生物学. 14 (4): 277–291. arXiv:1404.3686. 文献标识码:2014AsBio..14..277A. DOI:10.1089/ast.2013.1129. PMC 3995117. PMID 24611714.
\item 凯利，彼得（2013年7月18日）。“在寒冷恒星周围，一个更温暖的宜居行星，冰会变暖而非变冷”。华盛顿大学新闻。检索日期：2023年2月12日
\item 谢尔兹，A. L.；比特兹，C. M.；梅多斯，V. S.；乔希，M. M.；罗宾逊，T. D.（2014）。“恒星亮度增加引发的光谱驱动行星去冰化”。天体物理学杂志.785(1)：L9。arXiv：1403.3695文献编码2014ApJ...785L...9Sdoi10.1088/2041-8205/785/1/L9S2CID118544889
\item 巴恩斯，R.；穆林斯，K.；戈德布拉特，C.；梅多斯，V. S.；卡斯廷，J. F.；赫勒，R.（2013）。“潮汐金星：通过潮汐加热引发气候灾难”. 天体生物学. 13 (3): 225–250. arXiv:1203.5104. Bibcode:2013AsBio..13..225B. doi:10.1089/ast.2012.0851. PMC 3612283. PMID 23537135.
\item 赫勒，R.；阿姆斯特朗，J.（2014）。“超宜居世界”。天体生物学. 14 (1): 50–66. arXiv:1401.2392. Bibcode:2014AsBio..14...50H. doi:10.1089/ast.2013.1088. PMID 24380533. S2CID 1824897.
\item 杰克逊，B.；巴恩斯，R.；格林伯格，R.（2008）。“类地系外行星的潮汐加热及其对宜居性的影响”. 《皇家天文学会月报》. 391 (1): 237–245. arXiv:0808.2770. 天文学文献标识码:2008MNRAS.391..237J. DOI:10.1111/j.1365-2966.2008.13868.x. S2CID 19930771.
\item 吉尔斯特，保罗；勒佩奇，安德鲁（2015年1月30日）。“宜居行星候选名单综述”。半人马座梦想，陶零基金会。检索日期：2015年7月24日。
\item “美国国家航空航天局开普勒望远镜首次发现位于‘宜居带’的类地行星”.美国国家航空航天局，2014年4月17日。
\item 比尼亚米，乔瓦尼·F.（2015）。七重天之谜：智人将如何征服太空。施普林格出版社。第110页。ISBN 978-3-319-17004-6.
\item 霍威尔，伊丽莎白（2013年2月6日）。“距离最近的‘外星地球’可能只有13光年”。太空网。科技媒体网络。检索日期：2013年2月7日
\item 科普拉普，拉维·库马尔（2013年3月）。“对开普勒M型矮星宜居带内类地行星出现率的修订估算”。天体物理学报快报. 767(1)：L8。arXiv:1303.2649. 文献编码:2013ApJ...767L...8K. doi:10.1088/2041-8205/767/1/L8. S2CID 119103101.
\item “美国国家航空航天局开普勒任务发现地球的更大、更古老的表亲”，2015年7月23日，检索日期2015年7月23日。
\item 埃姆斯帕克，杰西（2011年3月2日）。“开普勒发现奇特的行星系统”。国际商业时报。国际商业时报公司。检索日期：2011年3月2日
\item “2010年NAM在格拉斯哥大学”。于2011年7月16日从原文存档。检索日期：2010年4月15日。
\item 萨特，保罗·M.（2022年12月9日）。“交换空间：恒星互换如何形成热木星”。宇宙今日。
 information@eso.org。“HIRES | ELT | ESO”。elt.eso.org。检索日期：2025年12月18日
 穆图，克莱尔；德勒伊，玛格丽；吉约，特里斯坦；巴格林，安妮；博尔德，帕斯卡尔；布希，弗朗索瓦；卡布雷拉，胡安；奇兹马迪亚，西拉德；德格，汉斯·J.（2013年11月1日）。“CoRoT：系外行星计划的成果”. 《伊卡洛斯》. 226 (2): 1625–1634. arXiv:1306.0578. 文献标识码:2013Icar..226.1625M. DOI:10.1016/j.icarus.2013.03.022. ISSN 0019-1035.
 市长，M；等（2003）。“用HARPS设定新标准”（PDF）。《信使》| 欧洲南方天文台。第114卷，第20页。文献标识码:2003Msngr.114...20M。检索日期2025年12月17日
 “凌日系外行星巡天卫星（TESS）”. 系外行星探索：太阳系外的行星。加州理工学院。 Public Domain本文引用了来自该来源的文字，该来源属于公共领域.
 \end{enumerate}